\documentclass[preprint,12pt,a4paper]{elsarticle}
\usepackage{amsmath,amssymb,amsthm,mathtools,booktabs,array}
\usepackage[hidelinks]{hyperref}
\usepackage{microtype}
\hypersetup{pdftitle={Marcinkiewicz--Zygmund constants of signed Padua rules beyond half-exactness},
 pdfauthor={},pdfsubject={Padua points; Marcinkiewicz--Zygmund inequalities; signed cubature; hyperinterpolation}}
\newtheorem{theorem}{Theorem}[section]
\newtheorem{lemma}[theorem]{Lemma}
\newtheorem{proposition}[theorem]{Proposition}
\newtheorem{corollary}[theorem]{Corollary}
\theoremstyle{remark}
\newtheorem{remark}[theorem]{Remark}
\newcommand{\Pn}{\Pi}
\newcommand{\R}{\mathbb R}
\newcommand{\CC}{\mathrm{CC}}
\newcommand{\Pa}{\mathrm P}
\newcommand{\lamin}{\lambda_{\min}}
\newcommand{\lamax}{\lambda_{\max}}
\newcommand{\norm}[1]{\left\lVert#1\right\rVert}
\newcommand{\dd}{\,\mathrm d}
\DeclareMathOperator{\diag}{diag}
\date{13 November 2026}
\journal{Journal of Approximation Theory}
\biboptions{sort&compress}
\emergencystretch=2em % narrower elsarticle text block
% elsarticle's preprint footer (\ps@pprintTitle) contains a stray space that makes the
% footer box 2.6pt overfull on page 1; tolerate it inside the output routine only.
\makeatletter\pretocmd{\@outputpage}{\hfuzz=3pt }{}{}\makeatother
\begin{document}
\begin{frontmatter}
\title{Marcinkiewicz--Zygmund constants of signed Padua rules\texorpdfstring{\\}{ }beyond half-exactness}
\author{[Author names to be inserted]}
\begin{abstract}
The interpolatory cubature rule on the Padua points of degree $m$ integrates
exactly all polynomials of degree $m$ on the square, but its weights have both
signs. We study its Marcinkiewicz--Zygmund constants, the extreme eigenvalues of
its Gram form on polynomials of degree $n$, beyond the half-exactness range
$n\le m/2$, where exactness makes the Gram matrix the identity. The marginals
of the rule are the Clenshaw--Curtis rule, with a factor two, and this
bounds the extreme eigenvalues by the Clenshaw--Curtis ones. For every fixed
ratio $n/m>1/2$ the least eigenvalue tends to $-\infty$ and the greatest to
$+\infty$, so half-exactness is sharp. At $n=m-1$ both extremes are of order
$\log m$: their normalized limits are the extreme points of the spectrum of an
explicit compact operator whose kernel is built from complete elliptic
integrals. Its norm is the modulus of its least eigenvalue, which is simple, and
we enclose it by a computer-assisted certificate. In contrast, the direct signed
coefficient hyperinterpolant has a $C\to L^2$ norm below thirty, uniformly. Exact
spectral brackets at six degrees up to $80$ and certified negative
directions at three degrees above $1000$ complete the picture.
\end{abstract}

\begin{keyword}
Padua points \sep Marcinkiewicz--Zygmund inequalities \sep signed cubature \sep hyperinterpolation \sep Gram matrix
\MSC[2020] 41A10 \sep 41A63 \sep 65D32 \sep 41A17
\end{keyword}
\end{frontmatter}

\section{Introduction}\label{sec:intro}

A Marcinkiewicz--Zygmund (MZ) inequality compares a discrete quadrature form
with the continuous $L^2$ norm on a polynomial space. If $Q$ is a quadrature
rule for a measure $\mu$ and $G$ is its Gram matrix on $\mathcal P_n$ in an
$L^2(\mu)$-orthonormal basis, the extreme eigenvalues
$\lamin(G)$ and $\lamax(G)$ are the sharp constants in
$\lamin(G)\|p\|^2\le Q(p^2)\le\lamax(G)\|p\|^2$, and the defect
$\eta=\|I-G\|_2$ measures the distance of $G$ from the identity. A defect below
one is what stability arguments for quadrature-based Galerkin schemes
require: the theorem of \citet{Wu26}, formulated for positive-weight rules,
asks for a defect smaller than the ratio of the coercivity to the continuity
constant of the bilinear form. If $Q$ is exact
through degree $d$, then $G=I$ for all degrees up to $\lfloor d/2\rfloor$
(the half-exactness range). Beyond this range nothing is forced by exactness,
and the question is how the extreme eigenvalues behave as the degree $n$ of the
polynomial space increases towards the degree $m$ of the rule.

\paragraph{The signed Padua rule}
The Padua points of degree $m$ are the $\dim\mathcal P_m$ points of the set
$\mathcal X_m$ defined in \eqref{eq:padua}, a checkerboard subset of a Chebyshev--Lobatto
grid on the square. Lagrange
interpolation at the Padua points was analysed by \citet{BCDMVX06}, whose
Theorem~2 gives the fundamental polynomials by a corrected-kernel formula, and
\citet{Xu96} studied interpolation at parity-separated
Chebyshev--Lobatto points of this type. Integrating the interpolant gives the
cubature rule $Q_m^{\mathrm P}$ for Lebesgue measure, whose weights are
formula~(33) of \citet{SVZ08}. It is exact on $\mathcal P_m$, its weights have
both signs, and negative weights were noted there to be few and small
\citep[p.~421]{SVZ08}. The positive rule with Chebyshev weights at the same points
is a different object: the two-sided $L^p$ MZ inequalities of
\citet{KLP18} at Padua-type nodes are for positive weights in the
Chebyshev-weighted norm, with unspecified constants, and the hyperinterpolation
of \citet{CDMV07} on the square uses positive Chebyshev-measure weights
at Xu points. In the present setting, Lebesgue measure and signed weights, the
finite-degree computations of \citet{ASV27} for a Padua-based cubature rule for
Lebesgue measure, with algebraic degree of exactness $m\le20$, give $\eta<1$ up to
$n=m-1$. The same authors write that such behaviour beyond half-exactness is, to
their knowledge, not known theoretically for Clenshaw--Curtis rules, which enter
the present paper through the marginals of the Padua rule. Here $m$ is always the degree of the Padua rule;
all logarithms are natural.

\paragraph{Results}
Write $G^{\mathrm P}_{m,n}$ for the Gram matrix of $Q_m^{\mathrm P}$ on
$\mathcal P_n$ with respect to Lebesgue measure on $[-1,1]^2$ (total mass four),
and $\lambda^{\mathrm P}_\pm(m,n)$ for its extreme eigenvalues.
\begin{itemize}
\item[\textbf{A.}] \emph{Marginal and compression} (Theorem~\ref{thm:padua}).
 For every rule exact on $\mathcal P_m$, the weights summed over one coordinate
 equal twice the Clenshaw--Curtis weights, and $\lamin(G^{\mathrm P}_{m,n})$ is
 at most, and $\lamax(G^{\mathrm P}_{m,n})$ at least, the corresponding
 Clenshaw--Curtis extreme eigenvalue for $m$ intervals. In particular
 $\liminf_m\|I-G^{\mathrm P}_{m,m-1}\|_2\ge1$.
\item[\textbf{B.}] \emph{A sharp threshold} (Theorem~\ref{thm:padua-carriers},
 Proposition~\ref{prop:padua-odd-carriers} and
 Corollary~\ref{cor:padua-ratio}). For the actual signed rule and every fixed
 degree ratio $\tau\in(\tfrac12,1)$, $\lamin(G^{\mathrm P}_{m,\lfloor\tau m\rfloor})\to-\infty$
 and $\lamax(G^{\mathrm P}_{m,\lfloor\tau m\rfloor})\to+\infty$ as $m\to\infty$,
 while $G^{\mathrm P}_{m,n}=I$ for $n\le m/2$. The proof uses explicit
 carriers and the logarithmic boundary-weight asymptotics of
 Theorem~\ref{thm:padua-weights}. Consequently no uniform strong MZ constant
 below one exists for any fixed ratio above one half.
\item[\textbf{C.}] \emph{The corner operator} (Theorems~\ref{thm:log-reduction}
 and~\ref{thm:corner-operator}, Section~\ref{sec:structure}). At $n=m-1$
 the quantities $\lambda^{\mathrm P}_\pm(m,m-1)/\log m$ converge to
 constants $c_-<0<c_+$, the extremes of the spectrum of an explicit compact
 operator $B$. The operator has an elliptic-integral kernel
 (Proposition~\ref{prop:elliptic-kernel}); its least eigenvalue is simple and
 equals $-\|B\|$, and $0<c_+<\|B\|$ (Proposition~\ref{prop:positivity}). A
 computer-assisted certificate encloses $\|B\|$ in a coarse interval
 (Proposition~\ref{prop:enclosure}), and analytic carrier bounds give
 explicit lower bounds (Corollary~\ref{cor:norm-lower}). No decimal digit of $c_\pm$ is
 determined here.
\item[\textbf{D.}] \emph{A uniform norm} (Theorem~\ref{thm:padua-operator}). The direct signed
 coefficient operator $H_{m,n}f=\sum_iQ_m^{\mathrm P}(fP_i)P_i$ satisfies
 $\|H_{m,n}\|_{C\to L^2}<30$ for all $m\ge2$ and $0\le n<m$, although the
 Gram form diverges at the endpoint degree.
\end{itemize}
Section~\ref{sec:numerics} gives certified spectral brackets at six degrees
between $10$ and $80$ (Proposition~\ref{prop:padua-brackets}), where the Gram form is positive
definite with $\eta<0.83$ at $n=m-1$, and certified negative directions at three degrees
above $1000$ (Proposition~\ref{prop:padua-finite-negative}), together with a
subsection of floating-point measurements that are labelled as such and used
nowhere in a proof. Section~\ref{sec:open} lists open problems.

\paragraph{Clenshaw--Curtis inputs}
The Clenshaw--Curtis rule enters through its weights and through three spectral
facts: $\lambda_+(N,n)\le2-4/N+4/N^2$ for $0\le n<N$, $\lambda_-(N,N-1)\to0$, and
$\lambda_-(512,511)<1/4$ (Proposition~\ref{prop:cc-inputs}). These three facts are
proved in the companion paper \citep{MZ-CC}, and are used here without proof.
The remaining facts about the Clenshaw--Curtis weights that we use are proved in
Section~\ref{sec:cc-recap}. All other statements of this paper are proved here; the
long proofs are collected in the two appendices.

\paragraph{Computer-assisted statements}
Three statements rest on exact or outward-rounded machine computation:
Propositions~\ref{prop:enclosure}, \ref{prop:padua-brackets}
and~\ref{prop:padua-finite-negative}. Each is labelled computer-assisted, the reduction to a finite check is
proved in the text, and the programs and receipts are listed in the section on data
and code availability.

\section{Conventions and Clenshaw--Curtis inputs}\label{sec:cc-recap}
Let $Q(f)=\sum_j w_j f(x_j)$ approximate a finite positive measure $\mu$,
and let $G=(Q(P_iP_j))_{ij}$ in a real $L^2(\mu)$-orthonormal basis
of a polynomial space. The basis-independent quantity
\begin{equation}\label{eq:eta}
 \eta=\norm{I-G}_2
 =\max\{1-\lamin(G),\lamax(G)-1\}
\end{equation}
is the smallest nonnegative constant in
\[
 |Q(p^2)-\norm p_{L^2(\mu)}^2|\le\eta\norm p_{L^2(\mu)}^2.
\]
In particular, $\eta<1$ means that every Gram eigenvalue lies in $(0,2)$.
If $Q$ is exact through degree $d$, exactness alone gives $G=I$ through
degree $\lfloor d/2\rfloor$. We call this the half-exactness range.
Throughout the CC analysis,
\begin{equation}\label{eq:nodes}
 N\ge2,\qquad \theta_j=\frac{j\pi}{N},\quad x_j=\cos\theta_j,
 \quad 0\le j\le N,
 \qquad Q_N(f)=\sum_{j=0}^N w_j f(x_j),
\end{equation}
where $w_j=\int_{-1}^1\ell_j(x)\dd x$ and $\ell_j\in\Pn_N$ is
the cardinal polynomial at $x_j$. The measure is unnormalized Lebesgue
measure $\dd x$. Write $G_{N,n}$ for the Gram matrix on $\Pn_n$ and
$\lambda_\pm(N,n)$ for its extreme eigenvalues. Thus
\begin{equation}\label{eq:variational}
 \lambda_-(N,n)=\inf_{0\ne p\in\Pn_n}
 \frac{Q_N(p^2)}{\int_{-1}^1p^2},\qquad
 \lambda_+(N,n)=\sup_{0\ne p\in\Pn_n}
 \frac{Q_N(p^2)}{\int_{-1}^1p^2}.
\end{equation}
Let $\varepsilon^{(N)}_0=\varepsilon^{(N)}_N=1/2$ and
$\varepsilon^{(N)}_j=1$ for $0<j<N$ be the endpoint factors of the
Chebyshev--Gauss--Lobatto mesh with $N+1$ points, and put
$d_j=\min(j,N-j)$. For $\theta\in[0,\pi]$ let
\begin{equation}\label{eq:cc-bracket}
 W_N(\theta)=1-\sum_{r=1}^{\lfloor(N-1)/2\rfloor}
      \frac{2\cos(2r\theta)}{4r^2-1}
 -\boldsymbol1_{2\mid N}\frac{\cos(N\theta)}{N^2-1}.
\end{equation}
The only facts on the Clenshaw--Curtis weights used below are the
following; they are proved here in a few lines.

\begin{lemma}[Clenshaw--Curtis weights]\label{lem:weight-comparison}
Let $N\ge2$.
\begin{enumerate}
\item[(i)] $w_j=\dfrac{2\varepsilon^{(N)}_j}{N}\,W_N(\theta_j)$ for
 $0\le j\le N$; in particular $w_0=w_N=N^{-2}$ for odd $N$ and
 $w_0=w_N=(N^2-1)^{-1}$ for even $N$.
\item[(ii)] For $0\le\theta\le\pi$,
 $|W_N(\theta)-\tfrac\pi2\sin\theta|\le1/N$ for odd $N$ and
 $\le N/(N^2-1)$ for even $N$.
\item[(iii)] For $1\le j\le N-1$,
 $w_j\ge(\pi-\tfrac43)\sin\theta_j/N\ge\sin\theta_j/N$.
\item[(iv)] $w_j\ge(1+d_j)/N^2$ for $0\le j\le N$; in particular
 $w_j>0$.
\end{enumerate}
\end{lemma}
\begin{proof}
(i) In this proof write $\varepsilon_j$ for $\varepsilon^{(N)}_j$.
For an integer $s$ the geometric sum gives
$\sum_{j=0}^N\varepsilon_j\cos(sj\pi/N)
=\frac12\sum_{j=0}^{2N-1}\cos(sj\pi/N)$, which equals $N$ if $2N\mid s$
and $0$ otherwise. With the product formula, for $0\le k,l\le N$,
\[
 \sum_{j=0}^N\varepsilon_j\cos(k\theta_j)\cos(l\theta_j)
 =\frac{N\delta_{kl}}{2\varepsilon_k}.
\]
That is, the square matrix $\Phi=(\cos k\theta_j)_{j,k}$ satisfies
$\Phi^{\mathsf T}\diag(\varepsilon_j)\Phi=\frac N2\diag(\varepsilon_k)^{-1}$, so
$\frac2N\diag(\varepsilon_k)\Phi^{\mathsf T}\diag(\varepsilon_j)$ is a left
inverse of $\Phi$, hence also a right inverse. This says that
\[
 \ell_j(\cos\theta)=\frac{2\varepsilon_j}{N}\sum_{k=0}^N
 \varepsilon_k\cos(k\theta_j)\cos(k\theta).
\]
Since $\int_{-1}^1T_k(x)\dd x=2/(1-k^2)$ for even $k$ and $0$ for odd $k$,
integrating gives the formula for $w_j$. The telescoping identity
\begin{equation}\label{eq:telescope-cc}
 \sum_{r=1}^q\frac2{4r^2-1}=1-\frac1{2q+1}
\end{equation}
gives $W_N(0)=W_N(\pi)=1/N$ for odd $N$ (take $q=(N-1)/2$) and
$W_N(0)=W_N(\pi)=1/(N-1)-1/(N^2-1)=N/(N^2-1)$ for even $N$
(take $q=N/2-1$); this is the endpoint statement.

(ii) The integrals
$\int_0^\pi\sin\theta\cos(k\theta)\dd\theta$ equal $2/(1-k^2)$ for even $k$
and $0$ for odd $k$. These cosine coefficients of $\sin\theta$ are
absolutely summable, so the cosine series converges uniformly, and by
uniqueness of cosine coefficients
\[
 \tfrac\pi2\sin\theta=1-2\sum_{r\ge1}\frac{\cos(2r\theta)}{4r^2-1}
 \qquad(0\le\theta\le\pi).
\]
With $q=\lfloor(N-1)/2\rfloor$, subtracting from \eqref{eq:cc-bracket} gives
$W_N(\theta)-\frac\pi2\sin\theta
=2\sum_{r>q}\cos(2r\theta)/(4r^2-1)-\boldsymbol1_{2\mid N}\cos(N\theta)/(N^2-1)$.
For odd $N$, \eqref{eq:telescope-cc} bounds its modulus by $1/(2q+1)=1/N$.
For even $N$ the term $r=N/2$ of the sum is $2\cos(N\theta)/(N^2-1)$, so
the expression equals
$\cos(N\theta)/(N^2-1)+2\sum_{r>N/2}\cos(2r\theta)/(4r^2-1)$, of modulus at
most $1/(N^2-1)+1/(N+1)=N/(N^2-1)$.

(iii), (iv) Let $1\le j\le N-1$. By concavity,
$\sin\theta_j=\sin(\pi d_j/N)\ge2d_j/N\ge2/N$. Hence the error in (ii) is
at most $\frac12\sin\theta_j$ for odd $N$, and at most
$N/(N^2-1)\le4/(3N)\le\frac23\sin\theta_j$ for even $N$ (this needs
$N\ge2$). Therefore $W_N(\theta_j)\ge(\frac\pi2-\frac23)\sin\theta_j$, and
(i) gives the first inequality of (iii); the second holds since
$\pi-4/3>1$. Moreover, by (iii) and
$\sin\theta_j\ge2d_j/N$, $w_j\ge(\pi-\frac43)\,2d_j/N^2\ge2d_j/N^2
\ge(1+d_j)/N^2$ for interior $j$, because $d_j\ge1$. At $j=0,N$,
(i) gives $w_j\ge N^{-2}$.
\end{proof}

\begin{proposition}[CC spectral inputs; \citealp{MZ-CC}]\label{prop:cc-inputs}
For $0\le n<N$, $\lambda_+(N,n)\le2-4/N+4/N^2<2$.
Moreover $\lambda_-(N,N-1)\to0$ as $N\to\infty$, and
$\lambda_-(512,511)<1/4$, so $\norm{I-G_{512,511}}_2>3/4$.
\end{proposition}

\section{The marginal and spectral compression}\label{sec:marginal}

Let $\mathcal P_n$ denote bivariate polynomials of total degree at most $n$.
For a construction degree $m\ge2$ the Padua points are
\begin{equation}\label{eq:padua}
 \mathcal X_m=\left\{\left(\cos\frac{j\pi}{m},
                   \cos\frac{k\pi}{m+1}\right):
    0\le j\le m,\ 0\le k\le m+1,\ j+k\text{ even}\right\},
\end{equation}
a set of $(m+1)(m+2)/2=\dim\mathcal P_m$ points, the checkerboard subset of the
tensor Chebyshev--Lobatto grid with $m$ and $m+1$ intervals.
Suppose real weights $W_{jk}$ on $\mathcal X_m$ define a rule $Q_m^\Pa$
that is exact on $\mathcal P_m$ for unnormalized Lebesgue
measure $\dd x\dd y$ on the square. This assumption applies to the
interpolatory rule obtained by integrating its interpolation basis, or by
solving its exact moment equations. The following implication needs neither
positivity of the weights nor a separate existence or uniqueness theorem
for that construction. Its univariate comparison object is the
Clenshaw--Curtis rule of Section~\ref{sec:cc-recap} on the mesh with $N=m$
intervals: $Q_N$, its Gram matrix $G_{N,n}$, and its extreme eigenvalues
$\lambda_\pm(N,n)$, with $N=m$ in this section.

\begin{theorem}[Marginal and compression]\label{thm:padua}
% IDs: PD-T021
At each $x$-coordinate the total weight equals
\begin{equation}\label{eq:marginal}
 \sum_{k:j+k\,\mathrm{even}}W_{jk}=2w_j,\qquad0\le j\le m,
\end{equation}
where $w_j$ are the Clenshaw--Curtis weights of the mesh $N=m$.
Let $G^\Pa_{m,n}$ be the real symmetric Gram matrix of $Q_m^\Pa$ on
$\mathcal P_n$ in a Lebesgue-orthonormal basis. For $0\le n\le m$,
\begin{align}
 \lamin(G^\Pa_{m,n})&\le\lambda_-(N,n),&
 \lamax(G^\Pa_{m,n})&\ge\lambda_+(N,n),\label{eq:padua-extreme}\\
 \norm{I-G^\Pa_{m,n}}_2&\ge\norm{I-G_{N,n}}_2&&(N=m).\label{eq:padua-eta}
\end{align}
In particular,
$\liminf_{m\to\infty}\norm{I-G^\Pa_{m,m-1}}_2\ge1$,
and the Padua defect at $m=512,n=511$ exceeds $3/4$.
\end{theorem}
\begin{proof}
Write $a_j=\sum_kW_{jk}$. Exactness on univariate $q(x)\in\Pn_m$ gives
$\sum_j a_jq(x_j)=2\int_{-1}^1q=\sum_j2w_jq(x_j)$.
Taking $q=\ell_j$ proves \eqref{eq:marginal}.
Define $\iota q(x,y)=q(x)/\sqrt2$. Then
$\int_{[-1,1]^2}(\iota q)^2=\int_{-1}^1q^2$, so $\iota$ is an isometry from
$\Pn_n$ into $\mathcal P_n$. The exact marginal gives for $p,q\in\Pn_n$
\[
 Q_m^\Pa((\iota p)(\iota q))=\tfrac12\sum_j a_jp(x_j)q(x_j)=Q_N(pq).
\]
This step uses the marginal rather than cubature exactness for the product,
whose degree can exceed $m$. In orthonormal coordinates,
$\iota^*G^\Pa_{m,n}\iota=G_{N,n}$ and $\iota^*\iota=I$.
The symmetric variational principle proves \eqref{eq:padua-extreme},
and compression of $I-G^\Pa_{m,n}$ proves \eqref{eq:padua-eta}.
No sign condition on $W_{jk}$ was used.
The endpoint lower eigenvalue tends to zero. Therefore
$\norm{I-G^\Pa_{m,m-1}}_2\ge1-\lambda_-(N,N-1)$ gives the stated
liminf after taking limits on the right.
Proposition~\ref{prop:cc-inputs} gives the finite consequence.
\end{proof}

\section{Signed weights: density and logarithmic boundary asymptotics}\label{sec:weights}

We now use the actual interpolatory Lebesgue weights $w_z$ of $Q_m^\Pa$,
rather than an arbitrary rule exact on $\mathcal P_m$.
Put $\theta_j=j\pi/m$, $\phi_k=k\pi/(m+1)$, $x_j=\cos\theta_j$,
$y_k=\cos\phi_k$, and let
$\varepsilon^{(m)}_j$ and $\varepsilon^{(m+1)}_k$ be the endpoint factors of the two
Lobatto meshes (Section~\ref{sec:cc-recap}). The positive auxiliary weights
\[
 \rho_{jk}=\frac{2\varepsilon^{(m)}_j\varepsilon^{(m+1)}_k}{m(m+1)},\qquad j+k\text{ even},
\]
have mass one; they are distinct from the signed weights $w_{jk}$ of
$Q_m^\Pa$. Put $c_0=1$, $c_a=2$ for $a>0$, and
$h_a=\int_{-1}^1T_a(x)\dd x$. Thus $h_a=0$ for odd $a$ and
$h_a=2/(1-a^2)$ for even $a$.

The following density is, up to the orientation of the axes and the
normalization of the basis, formula (33) of Sommariva, Vianello and
Zanovello \cite{SVZ08} for the weights of the integrated interpolant, written
in the product Chebyshev basis; the underlying cardinal kernel is
Theorem~2 of Bos, Caliari, De~Marchi, Vianello and Xu \cite{BCDMVX06}.
The short proof below fixes the exceptional-mode bookkeeping used afterwards.

\begin{lemma}[Integrated cardinal density]\label{lem:density}
% IDs: PD-T101, PD-T111
The signed weights satisfy $w_z=\rho_z d_m(z)$, where
\begin{equation}\label{eq:padua-density}
 d_m(x,y)=\sum_{a+b\le m}c_ah_ac_bh_bT_a(x)T_b(y)-2h_mT_m(x).
\end{equation}
\end{lemma}
\begin{proof}
Finite cosine sums on the two Lobatto grids show that
$\sqrt{c_ac_b}T_a(x)T_b(y)$, $a+b\le m$, are orthogonal for $\rho$,
with squared norm one except that $(a,b)=(m,0)$ has squared norm two.
Indeed the checkerboard sum is the full tensor sum plus its alternating
sum. A nonzero alternating product would require combined frequencies
at least $m$ and $m+1$, hence total degree at least $2m+1$;
products under consideration have degree at most $2m$. In the full tensor
sum only the highest $x$ mode has the exceptional Lobatto norm.
Their discrete reproducing kernel consequently gives the cardinal function
at $z=(x_j,y_k)$ as
\[
 L_z(x,y)=\rho_z\left\{\sum_{a+b\le m}c_ac_b
 T_a(x)T_a(x_j)T_b(y)T_b(y_k)-T_m(x)T_m(x_j)\right\}.
\]
There are exactly $\dim\mathcal P_m$ orthogonal functions and nodes,
so this kernel is cardinal. Integrating it proves
\eqref{eq:padua-density}; the constant $y$ factor in the exceptional
term has integral two.
\end{proof}

\begin{theorem}[Boundary weight limit]\label{thm:padua-weights}
% IDs: PD-T295
Let $e_0=1/2$ and $e_j=1$ for $j\ge1$. For fixed nonnegative integers
$j,k$ with $j+k$ even,
\begin{equation}\label{eq:padua-weight-limit}
 \lim_{\substack{m\to\infty\\m\text{ even}}}
 \frac{m^4w_{jk}}{\log m}=-16e_je_k(-1)^j.
\end{equation}
Moreover, for every even $m\ge4$ and every retained node,
\begin{equation}\label{eq:padua-weight-majorant}
 |w_{jk}|\le\frac{4096}{m^4}
       \bigl\{(1+j)(1+k)+\log(m+1)\bigr\}.
\end{equation}
The indices in this majorant are distances from the corner $(1,1)$;
the fixed-index limit is not asserted uniformly across the grid.
\end{theorem}
\begin{proof}
Write $m=2R$, $t_0=2$, $t_a=-4/(4a^2-1)$ for $a\ge1$, and
\[
 S_R(t,s)=\sum_{a+b\le R}t_at_b\cos(2at)\cos(2bs).
\]
Equation~\eqref{eq:padua-density} gives
\begin{equation}\label{eq:padua-weight-exact}
 w_{jk}=\frac{2\varepsilon^{(m)}_j\varepsilon^{(m+1)}_k}{m(m+1)}
 \{S_R(\theta_j,\phi_k)-2h_m\cos(m\theta_j)\}.
\end{equation}
The absolutely convergent cosine series
$f(t)=\pi\sin t=\sum_{a\ge0}t_a\cos(2at)$ on $[0,\pi]$
follows by integrating its coefficients. Equality follows from uniqueness
of cosine coefficients, or from density of polynomials in $\cos t$.
Let $T_R(t)=\sum_{a>R}t_a\cos(2at)$ and
\[
 V_R(t,s)=\sum_{\substack{1\le a,b\le R\\a+b>R}}
               t_at_b\cos(2at)\cos(2bs).
\]
Absolute convergence gives
\begin{equation}\label{eq:padua-triangle}
 S_R=f(t)f(s)-f(t)T_R(s)-f(s)T_R(t)+T_R(t)T_R(s)-V_R(t,s).
\end{equation}
Telescoping yields $|T_R|\le2/(2R+1)\le R^{-1}$.
Since $|t_a|\le4/(3a^2)$, split the omitted triangle into $a\le R/2$,
$b\le R/2$, and $a,b>R/2$. In the first strip there are at most $a$
choices of $b$, all greater than $R/2$, so it contributes at most
$64H_{\lfloor R/2\rfloor}/(9R^2)$; the second contributes the same.
The last part contributes at most $256/(9R^2)$. Hence
\[
 |V_R|\le64(1+\log R)/R^2.
\]
Using $\sin\theta_j\le\pi j/m$, $\sin\phi_k\le\pi k/m$ in
\eqref{eq:padua-triangle}, the bracket in
\eqref{eq:padua-weight-exact} has absolute value at most
\[
 m^{-2}\{\pi^4jk+2\pi^2(j+k)+260+256\log R+16/3\}.
\]
Multiplication by $2\varepsilon^{(m)}_j\varepsilon^{(m+1)}_k/[m(m+1)]\le2/m^2$
proves \eqref{eq:padua-weight-majorant}.

For fixed $j,k$, the terms other than $V_R$ in
\eqref{eq:padua-triangle} are $O_{j,k}(R^{-2})$. To evaluate $V_R$,
fix $0<\delta<1/2$. In its strip $a\le\delta R$, $R-a<b\le R$,
\[
 R^2t_b\cos(2b\phi_k)=-(-1)^k+O_k(\delta)+o_R(1),
 \quad \cos(2a\theta_j)=1+O_j(\delta^2),
\]
uniformly, while
$\sum_{a\le\delta R}a t_a=-\log R+O_\delta(1)$.
Thus this strip contributes $(-1)^k+O_{j,k}(\delta)+o_R(1)$
to $R^2V_R/\log R$. The strip $b\le\delta R$ contributes
$(-1)^j$ in the same way. The middle has absolute value at most
$16/(9\delta^4R^2)$, hence gives zero after this normalization.
Taking $R\to\infty$ and then $\delta\downarrow0$ gives
\[
 R^2V_R/\log R\longrightarrow(-1)^j+(-1)^k.
\]
Substitution into \eqref{eq:padua-weight-exact}, using
$(-1)^j=(-1)^k$ on the checkerboard and
$\varepsilon^{(m)}_j\varepsilon^{(m+1)}_k=e_je_k$ eventually, proves
\eqref{eq:padua-weight-limit}.
\end{proof}

% Source: PD-T377; topics/padua-mz-cleanup-20261105/theory/odd-weight-asymptotic.md
\begin{proposition}[Odd construction degrees]\label{prop:padua-odd-weights}
The fixed-index limit \eqref{eq:padua-weight-limit} also holds as $m$
tends to infinity through odd integers. For every odd $m\ge5$, the
whole-grid bound \eqref{eq:padua-weight-majorant} holds with $8192$
in place of $4096$.
\end{proposition}
\begin{proof}
Put $R=(m-1)/2$. Formula~\eqref{eq:padua-density} still gives
\eqref{eq:padua-weight-exact} with this $R$, and the exceptional term
is zero since $h_m=0$. The decomposition
\eqref{eq:padua-triangle} and its bounds hold for every integer $R\ge2$,
independently of the relation between $m$ and $R$. Now $m/R\le5/2$,
so
\[
 m^2|S_R(\theta_j,\phi_k)|\le\pi^4jk+
 \frac{5\pi^2}2(j+k)+\frac{25}4+400(1+\log R).
\]
Multiplication by $2\varepsilon^{(m)}_j\varepsilon^{(m+1)}_k/[m(m+1)]\le2/m^2$
proves the asserted majorant. In the thin-corner limit,
$2R\theta_j=j\pi(m-1)/m\to j\pi$ and
$2R\phi_k=k\pi(m-1)/(m+1)\to k\pi$.
Thus the two strips again give $(-1)^j+(-1)^k$ after normalization
by $R^2/\log R$, with the same negligible middle. Since
$m/R\to2$, \eqref{eq:padua-weight-limit} follows. No coordinate
interchange in the cardinal density is used.
\end{proof}

\section{Two signed carriers and the exact ratio threshold}\label{sec:carriers}

Define the degree-$r$ Fej\'er polynomial by
\begin{equation}\label{eq:padua-fejer}
 F_r(\cos t)=\frac1{r+1}
  \left(\frac{\sin((r+1)t/2)}{\sin(t/2)}\right)^2
 =1+2\sum_{\ell=1}^r\left(1-\frac\ell{r+1}\right)\cos(\ell t).
\end{equation}
Its value at $t=0$ is $r+1$. All norms in this section use
unnormalized Lebesgue measure.

\begin{theorem}[Signed countersequences]\label{thm:padua-carriers}
% IDs: PD-N105, PD-N106
Fix $a,b>0$ and put $r=\lfloor am\rfloor$, $s=\lfloor bm\rfloor$.
Along even $m$, if $a+b<1/2$, the polynomials
\[
 p_m^-(x,y)=T_{m/2}(x)F_r(x)F_s(y)
\]
have degree at most $m-1$ eventually and satisfy
\begin{equation}\label{eq:padua-negative}
 Q_m^\Pa((p_m^-)^2)/\log m\longrightarrow-16ab/9.
\end{equation}
If $2a+b<1/2$, the polynomials
\[
 p_m^+(x,y)=\tfrac12\{T_{m/2-r}(x)-T_{m/2+r}(x)\}F_r(x)F_s(y)
\]
also have degree at most $m-1$ eventually and satisfy
\begin{equation}\label{eq:padua-positive}
 Q_m^\Pa((p_m^+)^2)/\log m\longrightarrow8ab/9.
\end{equation}
Both sequences are nonzero and have squared norm at most $C_0^2$,
where $C_0=(1+\pi^4)/2$. Consequently
\[
 \lamin(G^\Pa_{m,m-1})\longrightarrow-\infty,\qquad
 \lamax(G^\Pa_{m,m-1})\longrightarrow+\infty
 \quad(m\text{ even}).
\]
\end{theorem}
\begin{proof}
The degrees are $m/2+r+s$ and $m/2+2r+s$ respectively, so the
strict parameter inequalities give admissibility. The second carrier
factor at $x=\cos t$ equals $\sin(mt/2)\sin(rt)$ and has modulus
at most one, as does $T_{m/2}$. For $L=r+1$,
\[
 |F_r(\cos t)|\le\min\{L,\pi^2/(Lt^2)\}.
\]
Splitting its squared integral at $1/L$ and using $\sin t\le t$
gives $\|F_r\|_2^2\le(1+\pi^4)/2=C_0$. Product integration proves
the asserted norm bound; the nonzero leading coefficients prove
nonvanishing.

For a mesh denominator $M$ and endpoint halfweights $c_j$, elementary
cosine sums give, for every cosine polynomial $h$ of degree $<M$,
\begin{equation}\label{eq:padua-trapezoid}
 \sum_{j=0}^M c_jh(j\pi/M)=M\widehat h_0,\qquad
 \sum_{j=0}^M c_j(-1)^jh(j\pi/M)=0.
\end{equation}
The constant cosine coefficient of $F_r^2$ is
\[
 D_r=1+2\sum_{\ell=1}^r(1-\ell/(r+1))^2
     =1+\frac{r(2r+1)}{3(r+1)},\qquad D_r/r\to2/3.
\]
For the negative carrier, the $x$ profile
$U_{m,j}=\cos^2(j\pi/2)[F_r(\cos\theta_j)/m]^2$
vanishes at odd $j$ and has weighted sum $D_r/(2m)\to a/3$.
The $y$ profile $V_{m,k}=[F_s(\cos\phi_k)/m]^2$ has weighted sum
$(m+1)D_s/m^2\to2b/3$ and alternating sum zero. These identities use
$2r<m$, $2s<m+1$, which hold eventually.

For the positive carrier replace the $x$ profile by
\[
 U^+_{m,j}=\sin^2(j\pi/2)\sin^2(r\theta_j)
                      [F_r(\cos\theta_j)/m]^2.
\]
It vanishes for even $j$. The polynomial
$h_r(\cos t)=\sin^2(rt)F_r(\cos t)^2$ has degree $4r<m$ and constant
coefficient $D_r/2-1/[2(r+1)^2]$: the coefficient of $T_{2r}$ in
$F_r^2$ is $2/(r+1)^2$. By \eqref{eq:padua-trapezoid},
$\sum_j\varepsilon^{(m)}_jU^+_{m,j}\to a/6$.

Here each profile converges at fixed indices, and the geometric bound
for $F_r$ gives $|F_r(\cos\theta_j)|/m\le C_a(1+j)^{-2}$
on the entire mesh, with the analogous $y$ bound. Together with
\eqref{eq:padua-weight-majorant}, the normalized signed summand is
bounded by a constant times $(1+j)^{-3}(1+k)^{-3}$. Extend arrays by
zero outside their grids and apply dominated convergence. The finite
profile identities therefore also identify their infinite limiting
sums. For the negative carrier only even $j$ and even $k$ remain,
so \eqref{eq:padua-weight-limit} gives
$-16(a/3)(b/3)=-16ab/9$. For the positive carrier only odd $j$ and
odd $k$ remain, giving $16(a/6)(b/3)=8ab/9$.
This proves \eqref{eq:padua-negative}--\eqref{eq:padua-positive}.
The continuous norms are positive and bounded above by $C_0^2$;
the variational principle then gives the two divergences. No finite
threshold is inferred from these limits.
\end{proof}

% Source: PD-T378; topics/padua-mz-cleanup-20261105/theory/odd-carriers.md
\begin{proposition}[Odd-degree carriers]\label{prop:padua-odd-carriers}
For odd $m$ put $M=m+1$, $r=\lfloor am\rfloor$, $s=\lfloor bm\rfloor$.
Under the respective parameter conditions of
Theorem~\ref{thm:padua-carriers}, replace its polynomials by
\[
 \widetilde p_m^-(x,y)=T_{M/2}(y)F_r(y)F_s(x),
\]
\[
 \widetilde p_m^+(x,y)=\tfrac12
    \{T_{M/2-r}(y)-T_{M/2+r}(y)\}F_r(y)F_s(x).
\]
Their degrees are at most $m-1$ eventually, their squared norms are
positive and at most $C_0^2$, and their signed-energy limits divided
by $\log m$ are respectively $-16ab/9$ and $8ab/9$.
Consequently both spectral divergences of
Theorem~\ref{thm:padua-carriers} hold over all integers $m\to\infty$.
\end{proposition}
\begin{proof}
The even mesh denominator is now $M$ in the $y$ coordinate. The
carrier values there are exactly $\cos(k\pi/2)$ and
$\sin(k\pi/2)\sin(r\phi_k)$. Their degrees are
$M/2+r+s$ and $M/2+2r+s$; the extra half in $M/2=m/2+1/2$ is
absorbed by the strict parameter inequalities. Polynomiality,
nonvanishing and the norm bounds follow just as before.

On the $y$ mesh the weighted profile sums are
$MD_r/(2m^2)\to a/3$ for the negative carrier, supported on even $k$,
and $M\{D_r/2-1/[2(r+1)^2]\}/(2m^2)\to a/6$ for the positive
carrier, supported on odd $k$. The unmodulated $x$ envelope has
weighted sum $D_s/m\to2b/3$ and alternating sum zero, so each of
its parity sums tends to $b/3$. These identities follow from
\eqref{eq:padua-trapezoid}, since $2s<m$ and respectively $2r<M$
or $4r<M$ eventually. The same profile decay and
Proposition~\ref{prop:padua-odd-weights} justify dominated convergence
in the complete signed sums. Even-even samples give
$-16(b/3)(a/3)$, and odd-odd samples give $16(b/3)(a/6)$.
The positive bounded norm and variational principle give both
odd-degree divergences. Combining the two parity subsequences proves
the all-integer conclusions.
\end{proof}

% Sources: PD-N107 and PD-T378; topics/padua-mz-cleanup-20261105/theory/odd-carriers.md
\begin{corollary}[Sharp half-degree threshold]\label{cor:padua-ratio}
% IDs: PD-N107, PD-T378
For $0\le\tau\le1/2$ and $n\le\tau m$, exactness gives
$G^\Pa_{m,n}=I$. For every fixed $1/2<\tau<1$,
\[
 \begin{gathered}
 \lamin(G^\Pa_{m,\lfloor\tau m\rfloor})\to-\infty,\qquad
 \lamax(G^\Pa_{m,\lfloor\tau m\rfloor})\to+\infty\\
 \quad(m\to\infty\text{ through all integers}).
 \end{gathered}
\]
Thus there is no uniform strong signed MZ constant $\eta_\tau<1$
above half-exactness.
\end{corollary}
\begin{proof}
For the first statement $2n\le m$. For the second set
$a=b=(\tau-1/2)/4$ in the two carriers. Their even-$m$ degrees are
bounded respectively by $(1/2+2a)m<\tau m$ and
$(1/2+3a)m<\tau m$; on odd $m$ add $1/2$ to each bound.
This overhead is harmless eventually because both inequalities are strict.
The even-$m$ proof and its odd-$m$ extension therefore apply.
\end{proof}

\begin{remark}[Refuted conjectures and scope]\label{rem:refuted}
% IDs: PD-T028, PD-T110, PD-T116 (OPEN with binding REJECTED: refuted dispositions)
Theorem~\ref{thm:padua-carriers}, Proposition~\ref{prop:padua-odd-carriers}
and Corollary~\ref{cor:padua-ratio} show that, for the signed Padua rule, the
Gram form on $\mathcal P_{m-1}$ is not positive definite for all large $m$
(its least eigenvalue tends to $-\infty$), that its largest eigenvalue is
unbounded, and that no fixed ratio $n/m\in(\tfrac12,1)$ admits a uniform strong MZ constant.
Thus the natural conjectures that this form stays positive definite, that its
upper spectrum stays below a fixed threshold, that the defect at
$n=m-1$ converges to one, and that signed MZ inequalities with a strong
constant below one hold at every fixed degree ratio, are false for this rule.
These conclusions do not apply to the positive Chebyshev Padua cubature
or to tensor Clenshaw--Curtis rules.
\end{remark}

\begin{lemma}[Cosine transform of the Fej\'er profile]\label{lem:fejer-cosine}
% IDs: PD-T379
Put
\[
 g(u)=\Bigl(\frac{\sin(u/2)}{u/2}\Bigr)^4,\ g(0)=1,\qquad
 h(x)=x^2\log|x|,\ h(0)=0 .
\]
For real $q$ let
\begin{equation}\label{eq:delta4}
 \Delta^4h(q)=3h(q)-2h(q+1)-2h(q-1)
 +\tfrac12\{h(q+2)+h(q-2)\}
\end{equation}
(one half of the centered fourth difference of $h$ with unit step). Then
\[
 \int_0^\infty u\,g(u)\cos(qu)\dd u=\Delta^4h(q)\qquad(q\in\R).
\]
In particular $\int_0^\infty u\,g(u)\dd u=\Delta^4h(0)=4\log2$.
\end{lemma}
\begin{proof}
Since $\sin^4(u/2)=(3-4\cos u+\cos2u)/8$, we have
$u\,g(u)=2(3-4\cos u+\cos2u)/u^3$, which is integrable on $(0,\infty)$.
By product-to-sum, $u\,g(u)\cos(qu)=2H_q(u)/u^3$ with
$H_q(u)=\sum_tc_t\cos(tu)$, where $(t,c_t)$ runs through
$(q,3)$, $(q\pm1,-2)$, $(q\pm2,\frac12)$. Here $\sum_tc_t=0$ and
$\sum_tc_tt^2=0$, so $H_q(u)=O(u^4)$ and $H_q'(u)=O(u^3)$ as
$u\downarrow0$. Integrating by parts twice on $[\delta,X]$ and letting
$\delta\downarrow0$, $X\to\infty$ (the boundary terms $H_q/u^2$ and
$H_q'/u$ tend to zero at $0$ by these estimates and at $\infty$ since
$H_q,H_q'$ are bounded) gives
$\int_0^\infty2H_q(u)u^{-3}\dd u=\int_0^\infty H_q''(u)u^{-1}\dd u$.
Because $\sum_tc_tt^2=0$,
$H_q''(u)=-\sum_{t\ne0}c_tt^2[\cos(tu)-\cos u]$. For $t\ne0$ and
$\varrho>0$, differentiating in $t$ and integrating
$t/(\varrho^2+t^2)$ gives
$\int_0^\infty e^{-\varrho u}[\cos(tu)-\cos u]u^{-1}\dd u
=\frac12\log\frac{\varrho^2+1}{\varrho^2+t^2}$; letting
$\varrho\downarrow0$ (Dirichlet's test justifies the limit) yields
$\int_0^\infty[\cos(tu)-\cos u]u^{-1}\dd u=-\log|t|$. Hence the integral
equals $\sum_{t}c_tt^2\log|t|$, with zero frequencies contributing zero.
This is $\Delta^4h(q)$. At $q=0$, $h(0)=h(\pm1)=0$ and $h(\pm2)=4\log2$.
\end{proof}

% Source: PD-T379; topics/padua-mz-cleanup-20261105/theory/growth-carrier-bounds.md
\begin{proposition}[Normalized one-sided growth bounds]\label{prop:padua-growth}
% IDs: PD-T379
With $g$ as in Lemma~\ref{lem:fejer-cosine} put
\[
 \begin{gathered}
 J_-(a)=\int_0^\infty u g(u)\cos^2(u/(2a))\dd u,\\
 J_+(a)=\int_0^\infty u g(u)\sin^2(u/(2a))\sin^2u\dd u.
 \end{gathered}
\]
These are finite and positive, and
$J_-(a)=\tfrac12\{4\log2+\Delta^4h(1/a)\}$.
Let $a,b>0$ and let $p_m^-$, $p_m^+$ be the carriers of
Theorem~\ref{thm:padua-carriers} (even $m$) and
Proposition~\ref{prop:padua-odd-carriers} (odd $m$), for $a+b<1/2$
respectively $2a+b<1/2$. On each parity their squared continuous norms
tend to $4\log2\,J_-(a)$ and $4\log2\,J_+(a)$. Consequently, over all
integers $m$ and also separately on even and odd $m$,
\begin{equation}\label{eq:padua-growth-fixed}
 \begin{gathered}
 \limsup_{m\to\infty}\frac{\lamin(G^\Pa_{m,m-1})}{\log m}
 \le-\frac{16ab}{9\cdot4\log2\,J_-(a)}\quad(a+b<\tfrac12),\\
 \liminf_{m\to\infty}\frac{\lamax(G^\Pa_{m,m-1})}{\log m}
 \ge\frac{8ab}{9\cdot4\log2\,J_+(a)}\quad(2a+b<\tfrac12).
 \end{gathered}
\end{equation}
In particular, with
\[
 \kappa^{\mathrm{car}}_-=\frac1{9\cdot4\log2\,J_-(1/4)},\qquad
 \kappa^{\mathrm{car}}_+=\frac1{36\cdot4\log2\,J_+(1/8)},
\]
\begin{equation}\label{eq:padua-growth-bounds}
 \limsup_{m\to\infty}\frac{\lamin(G^\Pa_{m,m-1})}{\log m}
 \le-\kappa^{\mathrm{car}}_-,\qquad
 \liminf_{m\to\infty}\frac{\lamax(G^\Pa_{m,m-1})}{\log m}
 \ge\kappa^{\mathrm{car}}_+.
\end{equation}
No matching bound, sharp constant, or existence of an extremal spectral
limit is asserted by this proposition.
\end{proposition}
\begin{proof}
Let $L=r+1$ and substitute $u=Lt$ in the continuous envelope norm.
The exact identity
\[
 F_r(\cos(u/L))/L=
 \left\{\frac{\sin(u/2)}{L\sin(u/(2L))}\right\}^2
\]
gives a squared-norm integrand tending to $u g(u)$ and bounded by
$u\min(1,\allowbreak\pi^4/u^4)$, extended by zero beyond $L\pi$.
The carrier mesh denominator is $D_m=m$ for even $m$ and $D_m=m+1$
for odd $m$. As $D_m/L\to1/a$ and $r/L\to1$, its factors tend to
$\cos(u/(2a))$ and $\sin(u/(2a))\sin u$. They have modulus at most
one. Dominated convergence and product integration prove the norm
limits, with $\int_0^\infty ug=4\log2$ by Lemma~\ref{lem:fejer-cosine}.
The integrands are continuous and positive on nonempty
intervals, so their integrals are positive. Writing
$\cos^2(u/(2a))=\frac12\{1+\cos(u/a)\}$, the lemma gives the stated
expression for $J_-(a)$.

The signed-energy limits \eqref{eq:padua-negative}--\eqref{eq:padua-positive}
divided by these norm limits give carrier
Rayleigh coefficients $-16ab/[9\cdot4\log2\,J_-(a)]$ and
$8ab/[9\cdot4\log2\,J_+(a)]$.
The variational principle applies on each parity, which gives
\eqref{eq:padua-growth-fixed}. Fix $a=1/4$ and
let $b\uparrow1/4$ through admissible values for the negative
carrier; fix $a=1/8$ and let $b\uparrow1/4$ for the positive one.
Their denominators are independent of $b>0$, which gives
\eqref{eq:padua-growth-bounds}. No boundary-value packet is assumed
admissible.
\end{proof}

The raw energy suprema are $1/9$ and $1/36$, obtained on the closures
of $a+b<1/2$ and $2a+b<1/2$ respectively. They are not normalized
spectral constants: the nonunit continuous packet norms must be
divided out as above.

\begin{proposition}[Strict improvement of the negative carrier constant]\label{prop:carrier-improvement}
% IDs: PD-T533
Put
\[
 \kappa_{\mathrm{new}}=\frac{16\cdot\frac{51}{200}\cdot\frac{49}{200}}
   {9\cdot4\log2\;J_-(51/200)}.
\]
Over all integers $m\to\infty$,
\[
 \limsup_{m\to\infty}\frac{\lamin(G^\Pa_{m,m-1})}{\log m}
 \le-\kappa_{\mathrm{new}},\qquad
 \kappa_{\mathrm{new}}>\kappa^{\mathrm{car}}_-.
\]
The same bound holds separately along even and along odd $m$.
This proposition asserts only this one-sided bound; the existence of
the limit is Theorem~\ref{thm:corner-operator} below.
\end{proposition}
\begin{proof}
For $q>2$ the function $h(x)=x^2\log x$ is smooth on $[q-2,q+2]$, with
$h''''(x)=-2/x^2$. Four applications of the fundamental theorem of
calculus to the centered fourth difference \eqref{eq:delta4} give
\begin{align*}
 \Delta^4h(q)&=\tfrac12\int_{[0,1]^4}h''''(q-2+\sigma_s)\dd s
 =-\int_{[0,1]^4}(q-2+\sigma_s)^{-2}\dd s<0,\\
 \sigma_s&:=s_1+s_2+s_3+s_4 .
\end{align*}
All arguments are positive, so we may differentiate under the integral:
\[
 (\Delta^4h)'(q)=2\int_{[0,1]^4}(q-2+\sigma_s)^{-3}\dd s\ge2q^{-3},
\]
by Jensen's inequality for the convex function $x^{-3}$, since the
argument has mean $q$. Let $J(a)=J_-(a)=\frac12\{4\log2+\Delta^4h(1/a)\}$
(Proposition~\ref{prop:padua-growth}). For $a\in[\frac14,\frac{51}{200}]$ we
have $1/a>2$, and therefore
\[
 0<J(a)<2\log2<\tfrac75,\qquad
 J'(a)=-\frac{(\Delta^4h)'(1/a)}{2a^2}\le-a.
\]
(For $2\log2<7/5$: $\log2<7/10$ because $e^{7/10}>1+\frac7{10}+\frac{(7/10)^2}2
+\frac{(7/10)^3}6+\frac{(7/10)^4}{24}>2$.)
Put $\varpi(a)=a(\frac12-a)$. Then $\varpi'(a)=\frac12-2a\in[-\frac1{100},0]$
on this interval, and $\varpi(a)\,a=a^2(\frac12-a)\ge\frac1{16}\cdot\frac{49}{200}$.
Since $\varpi'\le0$ and $0<J<\frac75$,
\[
 \varpi'J-\varpi J'\ge-\frac1{100}\cdot\frac75+\frac{49}{3200}
 =\frac{21}{16000}>0 .
\]
Thus $\varpi/J$ is strictly increasing on $[\frac14,\frac{51}{200}]$. Now
$\kappa_{\mathrm{new}}=\frac{16}{9\cdot4\log2}\,\varpi(\frac{51}{200})/J(\frac{51}{200})$
and $\kappa^{\mathrm{car}}_-=\frac{16}{9\cdot4\log2}\,\varpi(\frac14)/J(\frac14)$
because $\varpi(\frac14)=\frac1{16}$; this proves
$\kappa_{\mathrm{new}}>\kappa^{\mathrm{car}}_-$.

For $a=\frac{51}{200}$ and every fixed $0<b<\frac{49}{200}$ we have
$a+b<\frac12$, and \eqref{eq:padua-growth-fixed} gives, on either parity
and hence over all integers,
$\limsup\lamin(G^\Pa_{m,m-1})/\log m\le-16ab/[9\cdot4\log2\,J(a)]$.
Letting $b\uparrow\frac{49}{200}$ in this scalar inequality proves the
stated bound. No negative-corner estimate is subtracted and no boundary
parameter family is assumed admissible.
\end{proof}

\begin{proposition}[A ceiling for the scalar carrier family]\label{prop:carrier-ceiling}
% IDs: PD-T537
For fixed $a,b>0$ with $a+b<\frac12$ put
$\kappa_-(a,b)=16ab/[9\cdot4\log2\,J_-(a)]$, the modulus of the
limiting normalized Rayleigh coefficient of the negative carrier in
\eqref{eq:padua-growth-fixed}. Then
\begin{equation}\label{eq:carrier-ceiling}
 \kappa_-(a,b)<U_-:=\frac1{9\cdot4\log2\,(16\log2-9\log3)}<\frac1{24}.
\end{equation}
In particular $\sup\kappa_-\le U_-$ over this entire scalar parameter
domain. This bounds only the limiting trial coefficients of the negative
carriers of Theorem~\ref{thm:padua-carriers} and
Proposition~\ref{prop:padua-odd-carriers} (either parity), including the
closure of the parameter domain; it is not asserted to be optimal or
attained, it does not bound the full minimum of the Gram form, linear
combinations, other envelopes, or the positive carriers, and it gives no
statement for parameters that depend on $m$.
\end{proposition}
\begin{proof}
By the proof of Proposition~\ref{prop:carrier-improvement},
$(\Delta^4h)'>0$ on $(2,\infty)$, and $\Delta^4h$ is continuous on
$[2,\infty)$. Hence $\Delta^4h(q)>\Delta^4h(2)$ for $q>2$. Since
$1/a>2$, we get $J_-(a)>J_{\mathrm{end}}:=\frac12\{4\log2+\Delta^4h(2)\}$.
With $h(2)=4\log2$, $h(3)=9\log3$, $h(4)=32\log2$, $h(1)=h(0)=0$,
\[
 \Delta^4h(2)=28\log2-18\log3,\qquad
 J_{\mathrm{end}}=16\log2-9\log3=\log(2^{16}/3^9)>0,
\]
because $2^{16}=65536>19683=3^9$. Moreover $ab<a(\frac12-a)\le\frac1{16}$, so
$16ab/9<\frac19$ and
$\kappa_-(a,b)<\{9\cdot4\log2\,J_{\mathrm{end}}\}^{-1}=U_-$.
For the last inequality in \eqref{eq:carrier-ceiling}: $e^2<8$ gives
$\log2>\frac23$, hence $4\log2>\frac83$; and $2^{16}>3^{10}$ gives
$J_{\mathrm{end}}>\log3>1$ (as $e<3$). Thus
$9\cdot4\log2\,J_{\mathrm{end}}>24$.
\end{proof}

For orientation only, $\kappa_{\mathrm{new}}\approx0.0296$ and
$U_-\approx0.0333$; no claim is made that either is the best value within
the family.

\section{The corner operator and logarithmic limits}\label{sec:corner}

The carriers of Section~\ref{sec:carriers} give one-sided bounds of order
$\log m$ for the extreme Gram eigenvalues. The mechanism behind them is
visible in Theorem~\ref{thm:padua-weights}: close to a corner of the square
the signed weights are of order $m^{-4}\log m$, with signs alternating along
the checkerboard and with the coefficient $-16e_je_k(-1)^j$, while the
weight majorant \eqref{eq:padua-weight-majorant} controls the whole grid.
Polynomials of degree $<m$ resolve the nodes within a bounded index distance
of a corner through a Bessel-type kernel. This suggests that, after division
by $\log m$, the signed Gram form is asymptotically a fixed operator built
from this kernel, so that the sharp constants are the extremes of its
spectrum. This section makes that precise: a uniform reduction of the
whole-grid signed form to an explicit logarithmic form
(Theorem~\ref{thm:log-reduction}), the limit of the Gram kernel at a corner
(Lemma~\ref{lem:corner-kernel}), and the compact limit operator with the
convergence of the normalized extreme eigenvalues
(Theorem~\ref{thm:corner-operator}). The proofs are in
Appendix~\ref{app:corner}. In Section~\ref{sec:structure} we describe the
operator itself: its elliptic-integral kernel, the simplicity of its lowest
eigenvalue, which equals minus its norm, and enclosures of that norm, one of
them computer-assisted (proofs in Appendix~\ref{app:structure}).

\subsection{The reduction and the limit operator}

Throughout this section $m\ge4$, $M=m+1$, $R=\lfloor m/2\rfloor$,
$\theta_j=j\pi/m$ $(0\le j\le m)$, $\phi_k=k\pi/M$ $(0\le k\le M)$,
$z_{jk}=(\cos\theta_j,\cos\phi_k)$, and sums over $(j,k)$ run over the pairs
with $j+k$ even, the nodes of $\mathcal X_m$. We write $\varepsilon^{(m)}_j$
and $\varepsilon^{(M)}_k$ for the endpoint factors of
Section~\ref{sec:cc-recap}, $d_j=\min(j,m-j)$ and $d'_k=\min(k,M-k)$.
Let $\mathcal P_{m-1}$ carry the norm $\|p\|^2=\int_{[-1,1]^2}p^2\dd x\dd y$,
so that $Q_m^\Pa(p^2)=\sum_zw_zp(z)^2$ and $\lambda^\Pa_\pm(m,m-1)$
(the extreme eigenvalues of $G^\Pa_{m,m-1}$) are the extreme values of
$Q_m^\Pa(p^2)/\|p\|^2$. On $\mathcal P_{m-1}$ define the quadratic form
\begin{equation}\label{eq:L-m}
 L_m[p]=-\frac2{mMR^2}\sum_{j+k\text{ even}}
   \varepsilon^{(m)}_j\varepsilon^{(M)}_k
   \bigl[\cos(2R\theta_j)+\cos(2R\phi_k)\bigr]\,p(z_{jk})^2 .
\end{equation}

\begin{theorem}[Whole-grid reduction to the logarithmic form]\label{thm:log-reduction}
% IDs: PD-T534
For every $m\ge4$ and every $p\in\mathcal P_{m-1}$,
\begin{align}
 \bigl|Q_m^\Pa(p^2)-\log R\;L_m[p]\bigr|&\le4096\,\|p\|^2,\label{eq:log-error}\\
 |L_m[p]|&\le120\,\|p\|^2.\label{eq:L-bound}
\end{align}
Let $\nu_-(m)$ and $\nu_+(m)$ be the infimum and supremum of
$L_m[p]/\|p\|^2$ over $0\ne p\in\mathcal P_{m-1}$. Then
$|\nu_\pm(m)|\le120$ and, for each of the two signs separately,
\begin{equation}\label{eq:log-spectral}
 \bigl|\lambda^\Pa_\pm(m,m-1)-\log R\,\nu_\pm(m)\bigr|\le4096,\qquad
 \Bigl|\frac{\lambda^\Pa_\pm(m,m-1)}{\log m}-\nu_\pm(m)\Bigr|
  \le\frac{4300}{\log m}.
\end{equation}
\end{theorem}

This is a reduction to an explicit finite family, valid for both parities
of $m$; it is not by itself convergence to a fixed limiting operator, and the
constant $4300$ is an asymptotic transfer bound only. The next two
statements supply the limit.

Let $\mathsf P_k$ be the Legendre polynomial with $\mathsf P_k(1)=1$, and
let $\mathcal K_{m-1}$ be the reproducing kernel of $\mathcal P_{m-1}$ in
$L^2([-1,1]^2,\dd x\dd y)$ (unnormalized measure of mass four):
\begin{equation}\label{eq:repro-kernel}
 \mathcal K_{m-1}\bigl((x,y),(x',y')\bigr)=\sum_{\substack{k,l\ge0\\k+l\le m-1}}
  (k+\tfrac12)(l+\tfrac12)\,\mathsf P_k(x)\mathsf P_k(x')\mathsf P_l(y)\mathsf P_l(y').
\end{equation}
Let $J_0(z)=\sum_{r\ge0}(-1)^r(z/2)^{2r}/(r!)^2$ be the Bessel function of
order zero, $\mathcal T=\{(a,b):a,b\ge0,\ a+b\le1\}$, and
\begin{equation}\label{eq:K-infty}
 \mathcal K_\infty\bigl((u,v),(u',v')\bigr)=\int_{\mathcal T}ab\,J_0(au)J_0(au')J_0(bv)J_0(bv')\dd a\dd b .
\end{equation}
Note $\mathcal K_\infty((0,0),(0,0))=\int_{\mathcal T}ab=1/24$.

\begin{lemma}[Corner kernel limit]\label{lem:corner-kernel}
% IDs: PD-T535
Let $m\ge1$ be an integer, $u,v,u',v'\ge0$, shifts
$\varsigma_u,\varsigma_v,\varsigma_{u'},\varsigma_{v'}\in[0,1]$ and corner signs
$\mathbf s=(s_1,s_2)$, $\mathbf s'=(s_1',s_2')\in\{-1,1\}^2$. Put
\begin{align*}
 z_m&=\Bigl(s_1\cos\frac u{m+\varsigma_u},\ s_2\cos\frac v{m+\varsigma_v}\Bigr),\\
 z_m'&=\Bigl(s_1'\cos\frac {u'}{m+\varsigma_{u'}},\ s_2'\cos\frac {v'}{m+\varsigma_{v'}}\Bigr),
\end{align*}
$S=u+v+u'+v'$,
$\Theta=\frac14(u^2+v^2+u'^2+v'^2)+\varsigma_uu+\varsigma_vv+\varsigma_{u'}u'+\varsigma_{v'}v'$.
Then
\begin{equation}\label{eq:kernel-limit}
 \Bigl|m^{-4}\mathcal K_{m-1}(z_m,z_m')-\boldsymbol1_{\mathbf s=\mathbf s'}\,
   \mathcal K_\infty\bigl((u,v),(u',v')\bigr)\Bigr|\le\frac{3+S+\Theta/4}m .
\end{equation}
In particular the convergence is uniform on every cube $[0,U]^4$, with the
constant $C_U=3+5U+U^2/4$ for arbitrary shifts, and $C_U=3+\frac92U+U^2/4$
for the shifts $\varsigma_u=\varsigma_{u'}=0$, $\varsigma_v=\varsigma_{v'}=1$ of the
Padua meshes. Equal corners have limit $\mathcal K_\infty$ and different
corners (including corners differing in both coordinates) have limit zero.
\end{lemma}

\begin{theorem}[The corner operator and the logarithmic limits]\label{thm:corner-operator}
% IDs: PD-T536
Let $\mathcal H=L^2(\mathcal T,ab\dd a\dd b)$ with
$\langle f,g\rangle=\int_{\mathcal T}fg\,ab\dd a\dd b$, let
\[
 \psi_{jk}(a,b)=J_0(\pi ja)J_0(\pi kb),\qquad e_0=\tfrac12,\ e_j=1\ (j\ge1),
\]
and, for integers $K\ge1$, define the finite-rank selfadjoint operator
\begin{equation}\label{eq:B-K}
 B_K=-16\sum_{\substack{(1+j)(1+k)\le K\\ j+k\text{ even}}}
      e_je_k(-1)^j\,\langle\cdot,\psi_{jk}\rangle\,\psi_{jk}.
\end{equation}
Then $B_K$ converges in operator norm to a compact selfadjoint operator
$B$ on $\mathcal H$, and
\begin{equation}\label{eq:B-tail}
 \|B-B_K\|\le\frac{120}K .
\end{equation}
With $c_-=\inf\sigma(B)$ and $c_+=\sup\sigma(B)$, the Gram extremes of the
exact signed Lebesgue Padua rule of mass four on $\mathcal P_{m-1}$ satisfy,
as $m\to\infty$ through all integers,
\begin{equation}\label{eq:corner-limits}
 \frac{\lambda^\Pa_-(m,m-1)}{\log m}\longrightarrow c_-<0,\qquad
 \frac{\lambda^\Pa_+(m,m-1)}{\log m}\longrightarrow c_+>0 .
\end{equation}
Both limits hold over all integers and hence separately along even and odd
$m$, with exactly the same constants. The operator $B$ and the constants
$c_\pm$ do not depend on the finite Padua problems. No decimal enclosure or
closed elementary evaluation of $c_\pm$ is asserted here.
\end{theorem}

The operator $B$ is the limit of the single-corner blocks; it is never used
for a Gram bound at finite $m$ in this paper. In \eqref{eq:B-K} the infinite
operator is defined through the truncations and \eqref{eq:B-tail}, not through
a rearrangement of a conditionally convergent series.

\begin{remark}[Finite comparison]\label{rem:finite-comparison}
Let $I_K=\{(r,s)\in\mathbb Z_{\ge0}^2:(1+r)(1+s)\le K\}$, $N_K=2|I_K|$,
$U=\pi(K-1)$, $C_U=3+5U+U^2/4$ and
\[
 E_K(m)=\frac{N_K(32+24\pi K)}{8m}+24N_K\sqrt{C_U/m}.
\]
For all integers $K\ge1$ and $m\ge\max(4,2K)$, and for both signs,
\begin{equation}\label{eq:finite-comparison}
 \Bigl|\frac{\lambda^\Pa_\pm(m,m-1)}{\log m}-c_\pm\Bigr|
 \le\frac{4300}{\log m}+\frac{240}K+E_K(m).
\end{equation}
These coarse errors prove convergence with an explicit rate for each $K$;
they do not certify the finite range in which the limit is a good
approximation, and no such claim is made. The proof is at the end of
Appendix~\ref{app:corner}.
\end{remark}

\begin{remark}[The two active corners]\label{rem:two-corners}
Only two of the four corners of the square carry a nonvanishing limit
coefficient: the corners $(\pm1,1)$ for even $m$ and $(1,\pm1)$ for odd $m$.
Their limiting Gram blocks coincide, because the kernel limit
\eqref{eq:K-infty} does not depend on the corner signs, and their
coefficients are the same. Hence, for each fixed $K$, the nonzero
eigenvalues of the finite-rank form obtained from \eqref{eq:L-m} by keeping
the nodes within $(1+d_j)(1+d'_k)\le K$ converge, counted with multiplicity,
to those of $B_K\oplus B_K$ (Lemma~\ref{lem:finite-rank}). The two active
corners therefore double the multiplicity of each limiting eigenvalue,
not its value, and $c_\pm$ are the extreme points of the spectrum of the
single-corner operator $B$. The factor $16=2\cdot8$ in \eqref{eq:B-K} already
contains the two strip contributions $\cos(2R\theta_j)$ and
$\cos(2R\phi_k)$ of \eqref{eq:L-m}; there is no further factor two from the
corners.
\end{remark}

\subsection{Structure of \texorpdfstring{$B$}{B} and enclosures of its norm}\label{sec:structure}

The operator $B$ of Theorem~\ref{thm:corner-operator} has an explicit kernel.
Let $\mathbf K(k)=\int_0^{\pi/2}(1-k^2\sin^2s)^{-1/2}\dd s$ ($0\le k<1$) be the
complete elliptic integral of the first kind. For $0<a,a'<1$ put
\begin{equation}\label{eq:sigma-kernels}
 \begin{aligned}
 \Sigma_0(a,a')&=\frac{2\,\mathbf K(k_0)}{\pi^2(a+a')},&
 k_0&=\frac{2\sqrt{aa'}}{a+a'},\\[1ex]
 \Sigma_1(a,a')&=\frac{2\,\mathbf K(k_1)}{\pi^2\sqrt{aa'}}\,\boldsymbol1_{a+a'>1},&
 k_1^2&=\frac{(a+a')^2-1}{4aa'},
 \end{aligned}
\end{equation}
which are defined off the null sets $a=a'$ and $a+a'=1$; for $a+a'>1$ one has
$0<k_1<1$. Formally, $\Sigma_t(a,a')=\sum_{n\ge0}e_n(-1)^{tn}J_0(\pi na)J_0(\pi na')$
for $t=0,1$, with $e_n$ as in Theorem~\ref{thm:corner-operator}. We write
$\dd\mu=ab\dd a\dd b$ for the measure of $\mathcal H$.

\begin{proposition}[Elliptic-integral kernel of $B$]\label{prop:elliptic-kernel}
% IDs: PD-T539
The functions $\Sigma_0,\Sigma_1$ are measurable, nonnegative and symmetric,
and for almost every $a\in(0,1)$
\begin{equation}\label{eq:row-bounds}
 \int_0^1\Sigma_0(a,a')\,a'\dd a'\le\frac1\pi,\qquad
 \int_0^1\Sigma_1(a,a')\,a'\dd a'\le\frac2\pi .
\end{equation}
The operator $B$ of Theorem~\ref{thm:corner-operator} is the bounded integral
operator on $\mathcal H=L^2(\mathcal T,\dd\mu)$,
\begin{equation}\label{eq:B-kernel}
 \begin{gathered}
 (Bf)(x)=\int_{\mathcal T}\beta(x,y)f(y)\dd\mu(y),\\
 \beta\bigl((a,b),(a',b')\bigr)=-8\bigl[\Sigma_1(a,a')\Sigma_0(b,b')
  +\Sigma_0(a,a')\Sigma_1(b,b')\bigr].
 \end{gathered}
\end{equation}
In particular the kernel $\beta$ is nonpositive almost everywhere. The
identification is that of the norm limit $B=\lim B_K$; no pointwise
meaning of the series \eqref{eq:B-K} (the formal series for $\Sigma_t$) is asserted.
\end{proposition}

Put $B^-=-B$, with the nonnegative kernel $-\beta$, and $r=\|B\|=\|B^-\|$.
The kernel $-\beta$ is not strictly positive: it vanishes where
$a+a'<1$ and $b+b'<1$. Its square is nevertheless positivity improving.

\begin{proposition}[Positivity, simplicity and the norm identity]\label{prop:positivity}
% IDs: PD-T540
With $B^-=-B$ and $r=\|B\|$, the following hold.
\begin{enumerate}
\item[(i)] For every $f\in\mathcal H$ with $f\ge0$ almost everywhere and
 $f\ne0$, $(B^-)^2f>0$ almost everywhere on $\mathcal T$.
\item[(ii)] $r>0$ and $c_-=\inf\sigma(B)=-r$.
\item[(iii)] The eigenvalue $c_-=-r$ of $B$ is simple, with a real
 eigenfunction $u>0$ almost everywhere (unique up to a scalar factor), and
 $r$ is not an eigenvalue of $B$. Together with $c_+>0$ from
 Theorem~\ref{thm:corner-operator},
 \[
  0<c_+<r=\|B\|=-c_- .
 \]
\end{enumerate}
\end{proposition}

The norm of $B$ is therefore its most negative eigenvalue, and
$c_+$ is strictly smaller in modulus. The next two statements enclose $r$.
For $0\le u<v\le1$, $t\in\{0,1\}$, $\chi_0=1$, $\chi_1=2$, $x_+=\max(x,0)$ and
$0\le a\le1$, let
\begin{equation}\label{eq:I-t}
 I_t(a;[u,v])=\frac{\chi_t}{\pi^2}\int_0^\pi
 \Bigl[\sqrt{\bigl(v^2-(t-a\cos\theta)^2\bigr)_+}
      -\sqrt{\bigl(u^2-(t-a\cos\theta)^2\bigr)_+}\Bigr]\dd\theta ,
\end{equation}
and for intervals $U,U'\subset[0,1]$ put $J_t(U,U')=\int_UI_t(a;U')\,a\dd a$ and,
for rectangles $U\times V$,
$P_{U\times V}(a,b)=8\bigl[I_1(a;U)I_0(b;V)+I_0(a;U)I_1(b;V)\bigr]$.
These are integrals of elementary functions of square roots; no elliptic
integral occurs in them.

\begin{proposition}[Rectangle certificates]\label{prop:rectangles}
% IDs: PD-T541
Let $0\le u<v\le1$ and $t\in\{0,1\}$.
\begin{enumerate}
\item[(a)] $I_t(\cdot\,;[u,v])$ is continuous with
 $0\le I_t(a;[u,v])\le\chi_t\sqrt{v^2-u^2}/\pi$, and for almost every
 $a\in(0,1)$, $I_t(a;[u,v])=\int_u^v\Sigma_t(a,a')\,a'\dd a'$.
 Consequently, for a rectangle $U\times V\subset[0,1]^2$,
 $P_{U\times V}(x)=\int_{U\times V}(-\beta)(x,y)\dd\mu(y)$ for almost every
 $x\in\mathcal T$.
\end{enumerate}
Partition $[0,1]^2$ by finitely many lines parallel to the axes, let
$\mathcal R_i=U_i\times V_i$ ($1\le i\le p$) be those of the resulting
rectangles that meet $\mathcal T$ in a set of positive measure, put
$C_i=\mathcal R_i\cap\mathcal T$, and let $h_i>0$.
\begin{enumerate}
\item[(b)] (upper certificate) If $Q_i\in\R$ satisfy
 $\sum_jh_jP_{\mathcal R_j}(x)\le Q_i$ for almost every $x\in C_i$ and all $i$,
 then $\|B\|\le\max_iQ_i/h_i$.
\item[(c)] (lower certificate) Let $F$ be the set of $i$ with
 $\mathcal R_i\subset\mathcal T$, and let
 $0\le\mathsf m_{ij}\le8[J_1(U_i,U_j)J_0(V_i,V_j)+J_0(U_i,U_j)J_1(V_i,V_j)]$ for $i,j\in F$.
 Then
 \[
  \|B\|\ \ge\ \frac{\sum_{i,j\in F}h_ih_j\mathsf m_{ij}}{\sum_ih_i^2\mu(C_i)} .
 \]
\end{enumerate}
\end{proposition}

\begin{proposition}[A certified enclosure of $\|B\|$; computer-assisted]\label{prop:enclosure}
% IDs: PD-T542
Let $L=n_L/2^{256}$ and $U=n_U/2^{256}$ with
\begin{align*}
 n_L&=15243966599763808152468357735730532561\\&\qquad56125250651221590007974186558840498736,\\
 n_U&=55326534787122342745740713698404372650\\&\qquad886854985644685910272386076268567928344 .
\end{align*}
Then $L\le\|B\|\le U$, equivalently $-U\le c_-\le-L$, and in decimal form
\[
 0.013164946500379018110831\ \le\ \|B\|\ \le\ 0.477809279990884889210594 .
\]
\end{proposition}

This is a computer-assisted statement. It is obtained by choosing $h_i\equiv1$
in Proposition~\ref{prop:rectangles}, on the $16\times16$ grid of equal
squares, and enclosing the finitely many integrals \eqref{eq:I-t} by exact
integer interval arithmetic with $256$ angular cells; the reduction and the
arithmetic are described in Appendix~\ref{app:certificate}, and the code and
the receipts are listed under data and code availability.
The enclosure is deliberately coarse and should not be read as a numerical
value: its endpoints differ by a factor larger than $36$, so it fixes no
decimal digit of $\|B\|$ or of $c_-$, and for $c_+$ it yields only
$0<c_+<0.4779$. No decimal digit of $c_-$ or $c_+$ is determined in this
paper. Sharper arithmetic (finer grids and better trial functions in
Proposition~\ref{prop:rectangles}) has not been carried out.

The analytic one-sided bounds of Section~\ref{sec:carriers} identify the same
quantities from below. Combining them with Theorem~\ref{thm:corner-operator}
and Proposition~\ref{prop:positivity} requires no new computation.

\begin{corollary}[Analytic lower bounds for $\|B\|$ and $c_+$]\label{cor:norm-lower}
% IDs: PD-T?mz-paper-claude-02
With $\kappa_{\mathrm{new}}$ of Proposition~\ref{prop:carrier-improvement} and
$\kappa^{\mathrm{car}}_+$ of Proposition~\ref{prop:padua-growth},
\[
 \|B\|=-c_-\ \ge\ \kappa_{\mathrm{new}},\qquad
 c_+\ \ge\ \kappa^{\mathrm{car}}_+>0 .
\]
\end{corollary}
\begin{proof}
By Theorem~\ref{thm:corner-operator}, $\lamin(G^\Pa_{m,m-1})/\log m\to c_-$.
By Proposition~\ref{prop:carrier-improvement} its upper limit is at most
$-\kappa_{\mathrm{new}}$, so $c_-\le-\kappa_{\mathrm{new}}$, and
Proposition~\ref{prop:positivity}(ii) gives $\|B\|=-c_-\ge\kappa_{\mathrm{new}}$.
Likewise $\lamax(G^\Pa_{m,m-1})/\log m\to c_+$ and, by
\eqref{eq:padua-growth-bounds}, its lower limit is at least
$\kappa^{\mathrm{car}}_+$; hence $c_+\ge\kappa^{\mathrm{car}}_+$.
\end{proof}

For orientation only (not part of the statements above),
$\kappa_{\mathrm{new}}\approx0.02964$ and $\kappa^{\mathrm{car}}_+\approx0.01278$.
Thus the closed-form lower bound for $\|B\|$ is about $2.25$ times the lower
endpoint $L$ of Proposition~\ref{prop:enclosure}; the upper endpoint $U$ is not
improved. The corollary gives no upper bound on $c_+$ beyond
$c_+<\|B\|\le U$, and no optimality statement for either carrier constant.
The proofs of Propositions~\ref{prop:elliptic-kernel}--\ref{prop:enclosure}
are in Appendix~\ref{app:structure}.

\section{A uniform signed hyperinterpolation norm}\label{sec:uniform-norm}

For a Lebesgue-orthonormal basis $P_1,\ldots,P_d$ of $\mathcal P_n$,
define the direct signed coefficient operator
\[
 H_{m,n}f=\sum_{i=1}^dQ_m^\Pa(fP_i)P_i.
\]
It is distinct from inverse-Gram least squares. Euclidean duality gives
\begin{equation}\label{eq:padua-L1-norm}
 \|H_{m,n}\|_{C\to L^2}
 =\sup_{\substack{p\in\mathcal P_n\\\|p\|_2=1}}
          \sum_z|w_z|\,|p(z)|.
\end{equation}
Indeed the inequality in one direction follows by duality and the
triangle inequality. Equality follows by assigning the signs of
$w_zp(z)$ as bounded continuous data at the finitely many distinct nodes.

\begin{theorem}[Uniform coefficient norm]\label{thm:padua-operator}
% IDs: PD-T297
For every $m\ge2$ and $0\le n<m$,
\[
 \begin{gathered}
 2\le\|H_{m,n}\|_{C\to L^2}\le C_{\Pa}<30,\\
 C_{\Pa}=\frac{4\sqrt2\,\pi}{\pi-4/3}
       +\frac{68}{3}\left(\frac12+\frac1\pi\right)
       =28.3763348311\ldots.
 \end{gathered}
\]
\end{theorem}
\begin{proof}
Let $g(x)=\pi\sqrt{1-x^2}$, $R=\lfloor m/2\rfloor$, and $M'=2R+1$.
The absolutely convergent coefficients $t_r$ used above have
$\sum|t_r|=4$. Since $\sum_{j=0}^q t_j=2/(2q+1)$, summing the triangle first in
its second coordinate gives
\[
 A_R=\frac4{M'}-8\sum_{i=1}^R
       \frac1{(2i-1)(2i+1)(M'-2i)}.
\]
Partial fractions, followed by the sums over $2i-1$, $2i+1$ and $M'-2i$,
give
\[
 A_R:=\sum_{r+s\le R}t_rt_s
 =\frac8{M'(M'+1)}-\frac{16}{M'^2-1}
              \sum_{i=1}^R\frac1{2i-1}\ge-\frac4{M'}.
\]
For the last inequality use
$\sum_{i=1}^R(2i-1)^{-1}\le(R+2)/3\le R(R+2)/M'$.
Changing signed coefficients to absolute values affects only the two
nonzero coordinate axes. Their total change is $16-16/M'$; hence
\[
 \sum_{r+s>R}|t_rt_s|=16/M'-A_R\le20/M'\le20/m.
\]
The exceptional term of \eqref{eq:padua-density} vanishes for odd $m$;
for even $m$, $2|h_m|=4/(m^2-1)\le8/(3m)$. Consequently
\begin{equation}\label{eq:padua-density-tail}
 \|d_m-g\otimes g\|_\infty\le68/(3m).
\end{equation}

Put $\dd\nu=\dd x/(\pi\sqrt{1-x^2})$ and $\kappa=1/2+1/\pi$.
For a univariate polynomial $q$ of degree $\le m-1$, the Legendre
expansion gives $\|q\|_\infty^2\le m^2\|q\|_2^2/2$.
In $\pi^{-1}\int_0^\pi q(\cos t)^2\dd t$, the middle interval
$[1/m,\pi-1/m]$ contributes at most $(m/2)\|q\|_2^2$ because
$\sin t\ge2/(\pi m)$; the two remaining intervals contribute at most
$(m/\pi)\|q\|_2^2$. Thus
\[
 \int q^2\dd\nu\le\kappa m\int q^2\dd x.
\]
Apply this in both coordinates. Discrete cosine orthogonality below
the exceptional mode gives
$\sum_z\rho_zp(z)^2=\int p^2\dd(\nu\otimes\nu)$
for $p\in\mathcal P_{m-1}$. Since $\sum\rho_z=1$,
\begin{equation}\label{eq:padua-raw-L1}
 \sum_z\rho_z|p(z)|\le\kappa m\|p\|_2.
\end{equation}

Let $v_z=\rho_zg(x_z)g(y_z)\ge0$, and put $\alpha=\pi-4/3$.
Lemma~\ref{lem:weight-comparison}\,(iii) gives
$w_{N,j}^{\CC}\ge\alpha\sin(j\pi/N)/N$ for $1\le j\le N-1$.
The bound has zero right side at the endpoints, where $w_{N,j}^{\CC}>0$ by
part~(iv) of the lemma. Hence
\[
 v_{jk}\le\frac{2\pi^2}{\alpha^2}w_{m,j}^{\CC}w_{m+1,k}^{\CC}.
\]
Enlarge the sum to the positive tensor grid and apply the bound
$\lambda_+<2$ of Proposition~\ref{prop:cc-inputs} twice, once in each
coordinate, to obtain
$\sum_zv_zp(z)^2\le8\pi^2\|p\|_2^2/\alpha^2$.
Reflection on the odd one of the two meshes cancels the alternating
sine sum, while
$\sum_{j=1}^{N-1}\sin(j\pi/N)=\cot(\pi/(2N))$.
It follows that
\[
 \sum_zv_z=\frac{\pi^2}{m(m+1)}
 \cot\frac\pi{2m}\cot\frac\pi{2(m+1)}\le4.
\]
Cauchy--Schwarz gives
$\sum_zv_z|p(z)|\le4\sqrt2\pi\|p\|_2/\alpha$.
Finally \eqref{eq:padua-density-tail} implies
$|w_z|\le v_z+[68/(3m)]\rho_z$. Combine this with
\eqref{eq:padua-raw-L1} and \eqref{eq:padua-L1-norm} for the upper
bound. Exactness reproduces the constant one, whose $L^2$ norm on
the square is two, proving the lower bound.
\end{proof}

\begin{remark}[Reproduction and error bound]\label{rem:reproduction}
% IDs: PD-T113
Uniformity of the norm does not require signed Gram positivity or an upper
signed MZ inequality. The operator reproduces low degrees: if
$q\in\mathcal P_k$ with $0\le k\le\min(n,m-n)$, then $H_{m,n}q=q$. Indeed
$qP_i$ has degree at most $k+n\le m$, where the rule is exact, so
$Q_m^\Pa(qP_i)=\int qP_i$, and $q=\sum_i(\int qP_i)P_i$ because $k\le n$.
Since $\|g\|_{L^2}\le2\|g\|_\infty$ on the square, Theorem~\ref{thm:padua-operator}
gives for continuous $f$ and $E_k(f)=\inf_{q\in\mathcal P_k}\|f-q\|_\infty$
\[
 \|f-H_{m,n}f\|_{L^2}\le\|f-q\|_{L^2}+\|H_{m,n}(f-q)\|_{L^2}
 \le(2+C_{\Pa})\|f-q\|_\infty ,
\]
hence $\|f-H_{m,n}f\|_{L^2}\le(2+C_{\Pa})E_k(f)$ for
$0\le k\le\min(n,m-n)$. At $n=m-1$ this only supplies $k\le1$, and this
cannot be improved (Proposition~\ref{prop:no-quadratics}). The uniform norm
alone is therefore not a fixed-function convergence theorem at that endpoint.
\end{remark}

\begin{proposition}[Quadratics are not reproduced at the endpoint]\label{prop:no-quadratics}
% IDs: PD-N051
For every $m\ge3$ the operator $H_{m,m-1}$ reproduces $\mathcal P_1$ but not
$\mathcal P_2$. More precisely, put $(q,p)=(T_2(x),T_{m-1}(x))$ for odd $m$ and
$(q,p)=(T_2(y),T_{m-1}(y))$ for even $m$; then $q,p\in\mathcal P_{m-1}$ and
\[
 \langle H_{m,m-1}q-q,\,p\rangle_{L^2}=
 \begin{cases}h_{m-1}-h_{m+1},& m\text{ odd},\\ h_m,&m\text{ even},\end{cases}
 \ \ne0 .
\]
Consequently no bound $\|f-H_{m,m-1}f\|_{L^2}\le C\,E_k(f)$ with $k\ge2$ and
$C$ independent of $f$ holds for any $m\ge3$.
\end{proposition}
\begin{proof}
Reproduction of $\mathcal P_1$ is Remark~\ref{rem:reproduction} with $k=1\le\min(m-1,1)$.
For $p\in\mathcal P_{m-1}$ and continuous $q$ the definition of $H_{m,m-1}$
gives $\langle H_{m,m-1}q,p\rangle=Q_m^\Pa(qp)$, so the left side equals
$Q_m^\Pa(qp)-\int_{[-1,1]^2}qp$. Recall $h_j=\int_{-1}^1T_j$, which is
$2/(1-j^2)$ for even $j$ and $0$ for odd $j$, and that a function of one
variable has the integral over the square equal to twice its integral over
$[-1,1]$.

Let $m\ge3$ be odd. Then $qp=\frac12(T_{m+1}(x)+T_{m-3}(x))$. At the nodes
$x_j=\cos(j\pi/m)$ one has $T_{m+1}(x_j)=T_{m-1}(x_j)=(-1)^j\cos(j\pi/m)$.
The rule is exact on $T_{m-1}(x)$ and $T_{m-3}(x)$, so
$Q_m^\Pa(qp)=h_{m-1}+h_{m-3}$, while $\int qp=h_{m+1}+h_{m-3}$. The
difference $h_{m-1}-h_{m+1}$ is nonzero because $m\pm1$ are distinct even
integers.

Let $m\ge4$ be even. Then $qp=\frac12(T_{m+1}(y)+T_{m-3}(y))$. At a node
$(x_j,y_k)$ with $j+k$ even, $T_{m+1}(y_k)=\cos k\pi=(-1)^k=(-1)^j=T_m(x_j)$,
so $Q_m^\Pa(T_{m+1}(y))=Q_m^\Pa(T_m(x))=2h_m$ by exactness. The term
$T_{m-3}(y)$ is integrated exactly, and $\int T_{m+1}(y)=0$ since $m+1$ is
odd. Hence the defect is $h_m=2/(1-m^2)\ne0$.

Finally $q\in\mathcal P_2$ gives $E_k(q)=0$ for $k\ge2$, whereas
$\|q-H_{m,m-1}q\|_{L^2}\ge|\langle H_{m,m-1}q-q,p\rangle|/\|p\|_{L^2}>0$.
\end{proof}

\section{Finite results and measurements}\label{sec:numerics}

This section collects finite statements. Propositions~\ref{prop:padua-brackets}
and~\ref{prop:padua-finite-negative} are computer-assisted theorems in exact
or outward-rounded arithmetic. Everything in Section~\ref{sec:measured} is a
floating-point measurement, labelled as such; none of it is used in a proof.

\subsection{Certified brackets at moderate degrees}

\begin{proposition}[Certified spectral brackets]\label{prop:padua-brackets}
% IDs: PD-T107, PD-T298, PD-T299
For $m\in\{10,20,30,40,60,80\}$ and $n=m-1$, the extreme eigenvalues of
$G^\Pa_{m,n}$ lie in the following open intervals, where each decimal endpoint
denotes the exact terminating rational number:
\begin{center}
\begin{tabular}{rrrcc}\toprule
$m$ & $n$ & $\dim\mathcal P_n$ & $\lamin(G^\Pa_{m,n})$ & $\lamax(G^\Pa_{m,n})$\\\midrule
10 & 9  & 55   & $(0.356191,\,0.356192)$ & $(1.634856,\,1.634857)$\\
20 & 19 & 210  & $(0.272233,\,0.272234)$ & $(1.723313,\,1.723314)$\\
30 & 29 & 465  & $(0.236705,\,0.236706)$ & $(1.757929,\,1.757930)$\\
40 & 39 & 820  & $(0.214817,\,0.214818)$ & $(1.778088,\,1.778089)$\\
60 & 59 & 1830 & $(0.188593,\,0.188594)$ & $(1.802011,\,1.802012)$\\
80 & 79 & 3240 & $(0.172730,\,0.172731)$ & $(1.816497,\,1.816498)$\\\bottomrule
\end{tabular}
\end{center}
In particular, at these six degrees the Gram form is positive definite and
$\|I-G^\Pa_{m,m-1}\|_2<0.83$. Nothing is asserted at other degrees: the brackets
establish neither positivity between the listed degrees nor any
all-degree bound.
\end{proposition}

The brackets are compatible with the divergences of Theorem~\ref{thm:padua-carriers},
which are asymptotic and say nothing at these degrees. They are also
consistent with the finite observation of An, Sommariva and Vianello
\cite{ASV27} that the Padua-based rule has $\eta<1$ up to $n=m-1$ for $m\le20$,
and extend it at the six listed degrees.

\begin{proof}[Proof (computer-assisted)]
Work in the product Chebyshev basis $T_a(x)T_b(y)$, $a+b\le m-1$, of
$\mathcal P_{m-1}$. Let $H$ and $Q$ be the matrices of the continuous form
$\int p^2$ and of the quadrature form $Q_m^\Pa(p^2)$ in this basis (the
polarized forms $\int pq$ and $Q_m^\Pa(pq)$ give the matrix entries). Both are
rational matrices and are computed exactly. For basis indices $(a,b)$, $(c,d)$,
the product formula $T_aT_c=\frac12(T_{a+c}+T_{|a-c|})$ gives
$H_{(a,b),(c,d)}=\frac14\sum_{r,s}h_rh_s$ and
$Q_{(a,b),(c,d)}=\frac14\sum_{r,s}\varrho(r,s)$, over
$r\in\{a+c,|a-c|\}$, $s\in\{b+d,|b-d|\}$, with $\varrho(r,s)=Q_m^\Pa(T_r(x)T_s(y))$.
To evaluate $\varrho$, let $A$ be the distance from $r$ to the nearest multiple
of $2m$ and $B$ the distance from $s$ to the nearest multiple of $2(m+1)$, so
that $T_r(x_j)=T_A(x_j)$ and $T_s(y_k)=T_B(y_k)$ at all nodes. If $A+B\le m$ then
$\varrho(r,s)=h_Ah_B$ by exactness. Otherwise, at a node with $j+k$ even,
$T_{m-A}(x_j)T_{m+1-B}(y_k)=(-1)^{j+k}T_A(x_j)T_B(y_k)=T_A(x_j)T_B(y_k)$ and
$(m-A)+(m+1-B)\le m$, so $\varrho(r,s)=h_{m-A}h_{m+1-B}$.
Since $H$ is positive definite, the eigenvalues of $G^\Pa_{m,m-1}$ are the
generalized eigenvalues of $(Q,H)$, and, for a rational $t$, the number of
eigenvalues below $t$ is the number of negative eigenvalues of $Q-tH$ (congruence
by $H^{-1/2}$).

To count them exactly, let $d_t>0$ be an integer with $A_t=d_t(Q-tH)$ integral,
and let $R_t$ be an integer matrix, obtained by rounding scaled approximate
eigenvectors computed in floating point (no accuracy of this computation is
used). The integer matrix $B_t=R_t^{\mathsf T}A_tR_t$ is formed exactly and
checked to be symmetric and strictly diagonally dominant, $|(B_t)_{ii}|>\sum_{j\ne i}|(B_t)_{ij}|$.
Then every matrix obtained by multiplying the off-diagonal entries by
$s\in[0,1]$ is nonsingular (at a maximal component of a null vector the row
inequality fails), so no eigenvalue crosses zero along this path and $B_t$ has
the inertia of its diagonal. In particular $B_t$ is nonsingular, so $R_t$ is
invertible, and by Sylvester's law $Q-tH$ has the same inertia. The
checks used are: at the two endpoints of each minimum bracket the numbers of
negative eigenvalues are $0$ and $1$, and at the two endpoints of each maximum
bracket they are $\dim\mathcal P_{m-1}-1$ and $\dim\mathcal P_{m-1}$. These counts
place exactly one eigenvalue in each endpoint interval, with no eigenvalue equal
to an endpoint. For even $m$ the matrices split into two blocks by the parity of the
$x$-degree, and the code works blockwise; the counts add. The checks were passed for every row, by two
separately written programs (see the data and code availability statement). The statement rests on the
exactness of the integer and rational arithmetic of the programs and on their
correctness.
\end{proof}

\subsection{Finite negative witnesses}

% Source: PD-T376; topics/padua-mz-cleanup-20261105/theory/finite-indefiniteness.md
\begin{proposition}[Finite negative witnesses; computer-assisted]\label{prop:padua-finite-negative}
% IDs: PD-T376
For each row there is a nonzero rational-coefficient polynomial
$p\in\mathcal P_{m-1}$ whose signed Rayleigh quotient belongs to the
displayed strict rational interval:
\begin{center}
\begin{tabular}{rrc}\toprule
$m$ & $n$ & $Q_m^\Pa(p^2)/\int_{[-1,1]^2}p^2$\\\midrule
1101 & 1100 & $(-35067/10^6,-35066/10^6)$\\
1240 & 1239 & $(-273/10^7,-270/10^7)$\\
1600 & 1599 & $(-23415/10^6,-23414/10^6)$\\\bottomrule
\end{tabular}
\end{center}
Consequently the signed Gram form on $\mathcal P_{m-1}$ is indefinite
at each displayed construction degree.
\end{proposition}
The displayed intervals bound trial Rayleigh quotients, not minimum
eigenvalues. No first sign change is identified: the proposition says nothing
about positivity at smaller degrees. The witness data are archived with the deposit
(see the data and code availability statement; DOI to be added).

\begin{proof}[Certificate reduction and trust]
The witness is
\[
 p=2^{-56}\sum_{a+b\le m-1}C_{ab}\,\mathsf P_a(x)\mathsf P_b(y)
\]
with integers $C_{ab}$ and Legendre polynomials normalized by $\mathsf P_a(1)=1$, hence an exact
rational polynomial. Its exact continuous squared norm
$N=\sum4C_{ab}^2/[2^{112}(2a+1)(2b+1)]$ is computed in exact rational arithmetic
from Legendre orthogonality and is positive; for $m=1600$ it is enclosed near
$1.00000000000000031419$. The weight density is exactly
\eqref{eq:padua-density}, including its exceptional mode. Two separately written
programs enclose the signed energy $Q_m^\Pa(p^2)$ and obtain a strictly negative
enclosure. The first uses outward-rounded ball arithmetic for the node tables
$\mathsf P_a(\cos(j\pi/m))$, $\mathsf P_b(\cos(k\pi/(m+1)))$ and the cosine tables of the weights,
and exact integer and rational arithmetic for the fixed-point tensor
contractions, with explicit table-error and rounding bounds. For $m=1600$ its
energy enclosure is $[-0.0234144\pm4.95\cdot10^{-8}]$. The second builds
outward node tables in binary64 and evaluates the contractions in floating
point with analytic forward-error bounds; this path assumes the IEEE binary64
rounding model and the standard dot-product error bound. The other two rows
are treated in the same way. No eigensolver accuracy is a premise: a negative
quotient exhibits a negative direction of the form, and the constant polynomial
has positive energy $Q_m^\Pa(1)=4$, so the form is indefinite. The intervals in the
table contain the quotients of the energy enclosures by $N$.
Of the even degrees for which a witness was certified, $m=1240$ is the smallest;
this does not decide positivity at any smaller degree.
\end{proof}

\subsection{Measurements}\label{sec:measured}

All numbers in this subsection are binary64 measurements (NumPy and SciPy
eigensolvers, with a matrix-free Gram operator at large $m$); they are
not enclosures and no proof uses them.
\begin{enumerate}
\item[(1)] In double precision the weights $w_z$ have exactly two negative entries at each of
 $m=10,20,30,40$; negative weights that are few and small were already noted by
 Sommariva, Vianello and Zanovello \cite{SVZ08}. The exact marginal identity
 \eqref{eq:marginal} holds to within $8.73\cdot10^{-16}$ at $m=10,20,31$.
\item[(2)] The minimum of the Gram spectrum on $\mathcal P_{m-1}$, at the
 following degrees (at $m=1600$ it lies in the interval of
 Proposition~\ref{prop:padua-finite-negative}):
\begin{center}
\begin{tabular}{rcc}\toprule
$m$ & $\lamin(G^\Pa_{m,m-1})$ & weight of the minimizer in mode~2\\\midrule
1600 & $-0.0234144$ & $0.9590$\\
1601 & $-0.0695784$ & $0.9635$\\
3000 & $-0.0813226$ & $0.9645$\\
3001 & $-0.1276739$ & $0.9680$\\\bottomrule
\end{tabular}
\end{center}
 Mode~2 is the span of the second eigenvalue of the finite form $L_m$ of
 \eqref{eq:L-m}, a doublet by Remark~\ref{rem:two-corners}.
\item[(3)] Least-squares fits $\alpha+\beta\log m$ of the minima over $m$ between
 $800$ and $3001$, by parity class and window, have slopes $\beta\approx-0.092$.
 Finite-window slopes are not the limits $c_\pm$: Theorem~\ref{thm:corner-operator}
 gives a limit of $\lambda^\Pa_\pm/\log m$ with no rate, and the slope of a
 window need not approximate $c_-$.
\item[(4)] Truncating $B$ at $K=128$ and evaluating the triangle integrals by Gauss
 rules of orders $128$ and $192$, the extreme eigenvalues of the finite-rank
 operator $B_K$ are about $-0.1336$ and $+0.0456$, and its second negative
 eigenvalue is about $-0.0955$. These proxies are not enclosures: the proved
 tail bound $\|B-B_K\|\le120/K$ equals $0.94$ at $K=128$ and certifies nothing
 at this precision.
\end{enumerate}
The secants of the minima between $m=1600$ and $3000$ and between $1601$ and
$3001$ have slopes $-0.0921$ and $-0.0925$ per unit of $\log m$, within $0.0034$ and $0.0031$ of the
second proxy in (4). This closeness was noticed a posteriori: the
criterion for it was not fixed before the data were seen (only the mode weight,
the distance to the second doublet and the decay rate of the compressed
branches were prescribed in advance), so it is an observation and not a
test of the mechanism. In particular these measurements do not locate $c_-$ or $c_+$.

\section{Open problems}\label{sec:open}

The following questions remain open for the signed Padua rule.
\begin{enumerate}
\item \emph{Decimal values of $c_\pm$.} By Corollary~\ref{cor:norm-lower}
 and Proposition~\ref{prop:enclosure},
 $\kappa_{\mathrm{new}}\le\|B\|=-c_-\le0.4779$ and
 $\kappa^{\mathrm{car}}_+\le c_+<\|B\|$; no decimal digit of either constant
 is determined. A two-sided enclosure of $c_-=-\|B\|$ and of $c_+$ that fixes
 digits would follow from finer instances of the certificates of
 Proposition~\ref{prop:rectangles} together with an upper certificate for $c_+$
 that is not available here. Whether $c_\pm$ have closed forms is unknown.
\item \emph{The first sign change.} The degree $m$ at which the form
 $Q_m^\Pa(p^2)$ on $\mathcal P_{m-1}$ first becomes indefinite is unknown:
 Proposition~\ref{prop:padua-brackets} certifies positive definiteness
 only at the six listed degrees, and Proposition~\ref{prop:padua-finite-negative}
 certifies indefiniteness only at $m=1101,1240,1600$. Whether the two parities
 cross at different degrees, and whether the form stays indefinite afterwards,
 is unknown.
\item \emph{Effective onsets and rates.} The divergences of
 Theorem~\ref{thm:padua-carriers} and Corollary~\ref{cor:padua-ratio} are
 asymptotic and give no explicit onset degree. Likewise
 Theorem~\ref{thm:corner-operator} gives no usable rate for
 $\lambda^\Pa_\pm(m,m-1)/\log m\to c_\pm$: the finite comparison of
 Remark~\ref{rem:finite-comparison} is explicit but too coarse to certify any
 finite range.
\item \emph{The sharp norm of the direct operator.} Theorem~\ref{thm:padua-operator}
 gives $2\le\|H_{m,n}\|_{C\to L^2}\le28.38$. The true supremum over $m$ and $n$,
 and the limit as $m\to\infty$ at $n=m-1$ (if it exists), are unknown.
\item \emph{Fixed-function convergence at the endpoint.} Whether
 $H_{m,m-1}f\to f$ in $L^2$ for each fixed continuous (or smooth) $f$ is open.
 The uniform norm bound and the reproduction of $\mathcal P_1$ alone do not give
 it, and quadratics are not reproduced
 (Proposition~\ref{prop:no-quadratics}).
\end{enumerate}
Positive Padua cubature for the Chebyshev measure, and tensor
Clenshaw--Curtis rules on the square, are not covered by any statement in this paper.

\section*{Data and code availability}\label{sec:data}
\begingroup\raggedright
% Commit hash and Zenodo DOI: placeholders for AC (no deposit exists yet).
Code and evidence are in the public repository
\url{https://github.com/aacchhee/padua_research} at commit
\texttt{[to be fixed at submission]}, archived as \texttt{[Zenodo DOI to be added]}.
Paths are relative to the repository root. All results not listed below are
proved analytically in this paper.

\emph{Proposition~\ref{prop:padua-brackets}} (certified brackets). Exact pencils,
proposed integer transforms and the exact diagonal-dominance checks are in
\path{topics/padua-mz-2d-frontier-20261105/code/cert_native.py} with the
C source \path{cert_pencil.c} in the same directory. A second program,
\path{independent_cert.py} (same directory), reconstructs the pencil in Python
rational arithmetic, uses fresh transforms and repeats the checks. The receipts
\path{evidence/cert-m10.json}, \dots, \path{evidence/cert-m80.json} and
\path{evidence/independent-cert-m10.json}, \dots,
\path{evidence/independent-cert-m80.json} in the same topic directory cover
$m=10,20,30,40,60,80$. For $m\le40$ the exact-rational checker is also available
as \path{tracks/mz-constants/code/mz2d_exact.py} with
\path{tracks/mz-constants/code/mz2d_independent_check.py}.

\emph{Proposition~\ref{prop:padua-finite-negative}} (negative witnesses).
The producer is \path{topics/padua-mz-cleanup-20261105/code/rigorous_witness.py}
and the separate checker \path{independent_witness.py} (same directory); small
receipts \path{evidence/witness-producer-m1101.json}, \dots, \path{-m1600.json} and
\path{evidence/witness-independent-m1101.json}, \dots, \path{-m1600.json}, and the
payload identities in \path{evidence/witness-payload-identities.json}, are
committed. The witness arrays themselves are not in the repository: the
three exact integer coefficient arrays $C_{ab}$ for $m=1101,1240,1600$
(compressed JSON files of about $4.2$, $5.4$ and $8.7$ MB) and the exact
rational norm files are archived with the deposit.

\emph{Propositions~\ref{prop:rectangles} and~\ref{prop:enclosure}} (enclosure of
$\|B\|$). The checker is
\path{topics/padua-mz-sharp-constants-20261106/code/P1_rectangle.py}, run as
\texttt{P1\_rectangle.py --N 16 --theta 256}, with the arithmetic library
\path{code/interval_dyadic.py} (integers scaled by $2^{-256}$, exact integer
square roots, a Machin enclosure of $\pi$, no floating arithmetic). The
receipts \path{evidence/P1-rectangle-N16.json} (the instance used) and
\path{evidence/P1-rectangle-N8.json} (a coarser instance) record the exact
endpoint integers; \path{evidence/check-P1-independent-N16.json} is a second run
that reproduces them bit for bit. The run takes about two seconds.

\emph{Measurements} (Section~\ref{sec:measured}): item (1) by
\path{tracks/mz-constants/code/reproduce.py} (suites \texttt{padua} and
\texttt{marginal}); items (2) and the mode weights by
\path{topics/padua-mz-sharp-constants-20261106/code/P2_modes.py} and the
receipts \path{evidence/P2-m1600.json}, \path{-m1601.json}, \path{-m3000.json},
\path{-m3001.json}; the fits and the proxies in (3), (4) by
\path{code/S4_padua_fit.py}, \path{code/S3_corner_model.py} and the receipts
\path{evidence/S4-padua-fit.json}, \path{evidence/S3-corner-model.json}, all in
the same directory.

\endgroup

\appendix
\renewcommand{\thesection}{\Alph{section}}
\section{Proofs for the corner operator}\label{app:corner}

Throughout, the notation of Section~\ref{sec:corner} is in force. Let
$t_0=2$, $t_a=-4/(4a^2-1)$ for $a\ge1$ and
$S_R(t,s)=\sum_{a,b\ge0,\ a+b\le R}t_at_b\cos(2at)\cos(2bs)$, as in the proof
of Theorem~\ref{thm:padua-weights}. Recall from that proof that
$f(t)=\sum_{a\ge0}t_a\cos(2at)=\pi\sin t$ on $[0,\pi]$, that
$T_R(t)=\sum_{a>R}t_a\cos(2at)$ satisfies $|T_R|\le\sum_{a>R}|t_a|=2/(2R+1)\le1/R$,
and that the identity \eqref{eq:padua-triangle} holds for every integer
$R\ge2$.

\subsection{Proof of Theorem~\ref{thm:log-reduction}}\label{app:reduction}

\paragraph{Weights for both parities}
Since $h_a=0$ for odd $a$ and $c_{2a}h_{2a}=t_a$ for $a\ge0$, only even
indices contribute to \eqref{eq:padua-density}, and
$T_m(x_j)=\cos(m\theta_j)$. Hence for every $m\ge2$, with
$R=\lfloor m/2\rfloor$,
\begin{equation}\label{eq:weight-both}
 w_{jk}=\frac{2\varepsilon^{(m)}_j\varepsilon^{(M)}_k}{mM}
 \bigl\{S_R(\theta_j,\phi_k)-2h_m\cos(m\theta_j)\bigr\},
\end{equation}
where $h_m=2/(1-m^2)$ for even $m$ and $h_m=0$ for odd $m$.
In particular $R^2|2h_m|=m^2/(m^2-1)<2$ for even $m$, and $h_m=0$ for odd $m$.

\begin{lemma}[Triangular tail, uniformly on the square]\label{lem:triangular-tail}
For an integer $R\ge2$ and $t,s\in[0,\pi]$ put $\tau=\min(t,\pi-t)$,
$\sigma=\min(s,\pi-s)$ and
$V_R(t,s)=\sum_{1\le a,b\le R,\ a+b>R}t_at_b\cos(2at)\cos(2bs)$. Then
\[
 \bigl|R^2V_R(t,s)-\log R\,[\cos(2Rt)+\cos(2Rs)]\bigr|\le53+\tfrac83R(\tau+\sigma).
\]
\end{lemma}
\begin{proof}
The cosines are unchanged if $t,s$ are replaced by $\tau,\sigma$.
Let $L=\lfloor R/2\rfloor$, so $R/4\le L\le R/2$ for $R\ge2$. The index
pairs split into three disjoint classes: $a\le L$ (then $b>R-a\ge R/2$),
$b\le L$, and $a,b\ge L+1$; indeed $a,b\le L$ would give $a+b\le R$.

\emph{First strip} ($1\le a\le L$, $R-a<b\le R$; exactly $a$ values of $b$,
all $>R/2$). For such $b$, $b^2-\frac14\ge\frac3{16}R^2$ (as $R\ge2$) and
$R^2-b^2\le2Ra$. Hence
\[
 |R^2t_b+1|=\frac{R^2-b^2+1/4}{b^2-1/4}\le\frac{32a}{3R}+\frac4{3R^2}\le\frac{12a}R ,
\]
using $a\ge1$. Also, since $|\cos x-1|\le|x|$ and $|\cos x-\cos y|\le|x-y|$,
$|\cos(2a\tau)\cos(2b\sigma)-\cos(2R\sigma)|\le2a\tau+2(R-b)\sigma\le2a(\tau+\sigma)$.
Therefore
$|R^2t_b\cos(2a\tau)\cos(2b\sigma)+\cos(2R\sigma)|\le12a/R+2a(\tau+\sigma)$.
With $|t_a|\le4/(3a^2)$, the $R^2$-weighted sum of the first strip differs
from $-\cos(2R\sigma)\sum_{a=1}^La\,t_a$ by at most
\[
 \sum_{a=1}^La|t_a|\Bigl[\frac{12a}R+2a(\tau+\sigma)\Bigr]
 \le\frac43L\Bigl[\frac{12}R+2(\tau+\sigma)\Bigr]\le8+\frac43R(\tau+\sigma).
\]
Next $-a\,t_a=4a/(4a^2-1)=1/a+1/[a(4a^2-1)]$, and
$0\le\sum_{a\ge1}1/[a(4a^2-1)]\le\frac13\sum_{a\ge1}a^{-3}\le\frac12$
(integral test). With $0\le H_L-\log L\le1$ and
$\log2\le\log(R/L)\le\log4$ we get $|-\sum_{a\le L}at_a-\log R|<4$.
So the first strip differs from $\log R\cos(2Rs)$ by at most
$12+\frac43R(\tau+\sigma)$. Interchanging the roles of $a$ and $b$, the
second strip differs from $\log R\cos(2Rt)$ by the same bound.

\emph{Middle} ($a,b\ge L+1>R/2$): $|t_a|,|t_b|\le16/(3R^2)$ and there are at
most $R^2$ pairs, so $R^2$ times the middle part has modulus at most
$256/9<29$. Adding, $12+12+29=53$.
\end{proof}

\begin{lemma}[CC comparison on the Padua meshes]\label{lem:tensor-cc}
Let $m\ge4$ and let $\alpha_j$ ($0\le j\le m$) and $\beta_k$
($0\le k\le M$) be the Clenshaw--Curtis weights of the meshes with $m$ and $M$
intervals. Then $\alpha_j\ge(1+d_j)/m^2$, $\beta_k\ge(1+d'_k)/M^2$, and for every
$p\in\mathcal P_{m-1}$
\begin{equation}\label{eq:tensor-bound}
 \sum_{j=0}^m\sum_{k=0}^M\alpha_j\beta_k\,p(z_{jk})^2\le4\|p\|^2 .
\end{equation}
Moreover $mM/R^2\le15/2$.
\end{lemma}
\begin{proof}
The lower bounds are Lemma~\ref{lem:weight-comparison}\,(iv), which also gives
$\alpha_j,\beta_k>0$. Fix $j$. As a polynomial in $y$, $p(x_j,\cdot)$ has degree
$\le m-1<M$, so Proposition~\ref{prop:cc-inputs} gives
$\sum_k\beta_kp(x_j,y_k)^2\le2\int_{-1}^1p(x_j,y)^2\dd y$. Multiply by
$\alpha_j>0$ and sum over $j$. For each fixed $y$, $p(\cdot,y)$ has degree
$\le m-1<m$ in $x$, so again by Proposition~\ref{prop:cc-inputs}
$\sum_j\alpha_jp(x_j,y)^2\le2\int_{-1}^1p(x,y)^2\dd x$. Integrating in $y$
proves \eqref{eq:tensor-bound}. Finally, $mM/R^2=4+2/R\le5$ for $m=2R$ and
$mM/R^2=4+6/R+2/R^2\le15/2$ for $m=2R+1$, $R\ge2$.
\end{proof}

\begin{lemma}[Tail estimate]\label{lem:tail-estimate}
Let $m\ge4$, $K\ge1$, and let $\mathcal N$ be any set of nodes $z_{jk}$ with
$(1+d_j)(1+d'_k)\ge K$. Then for every $p\in\mathcal P_{m-1}$
\[
 \frac4{mMR^2}\sum_{z_{jk}\in\mathcal N}\varepsilon^{(m)}_j\varepsilon^{(M)}_k\,p(z_{jk})^2
 \le\frac{120}K\|p\|^2 .
\]
\end{lemma}
\begin{proof}
By Lemma~\ref{lem:tensor-cc}, $(1+d_j)(1+d'_k)\le m^2M^2\alpha_j\beta_k$, so
$1\le m^2M^2\alpha_j\beta_k/K$ on $\mathcal N$. With
$\varepsilon^{(m)}_j\varepsilon^{(M)}_k\le1$, the left side is at most
\[
 \frac{4mM}{KR^2}\sum_{j,k}\alpha_j\beta_kp(z_{jk})^2
 \le\frac{4mM}{KR^2}\cdot4\|p\|^2\le\frac{120}K\|p\|^2,
\]
by \eqref{eq:tensor-bound} and $mM/R^2\le15/2$.
\end{proof}

\paragraph{Proof of \eqref{eq:log-error}, \eqref{eq:L-bound}, \eqref{eq:log-spectral}}
Let $R\ge2$. Multiplying \eqref{eq:padua-triangle} by $R^2$ and using
Lemma~\ref{lem:triangular-tail}, $|f(t)|\le\pi\tau$, $|f(s)|\le\pi\sigma$ and
$|T_R|\le1/R$,
\begin{equation}\label{eq:S-remainder}
 \bigl|R^2S_R(t,s)+\log R[\cos(2Rt)+\cos(2Rs)]\bigr|
 \le\pi^2R^2\tau\sigma+(\pi+\tfrac83)R(\tau+\sigma)+54 .
\end{equation}
At a node $(t,s)=(\theta_j,\phi_k)$ we have $\tau=\pi d_j/m$, $\sigma=\pi d'_k/M$
and $R\le m/2<M/2$, so $R\tau\le\pi d_j/2$, $R\sigma\le\pi d'_k/2$, and the right
side of \eqref{eq:S-remainder} is at most
$\frac{\pi^4}4d_jd'_k+\frac{\pi(\pi+8/3)}2(d_j+d'_k)+54\le54(1+d_j)(1+d'_k)=:54D_{jk}$,
because $\pi^4/4<25$ and $\pi(\pi+8/3)/2<10$. Let
$\ell_{jk}$ be the coefficient of $p(z_{jk})^2$ in \eqref{eq:L-m}. By
\eqref{eq:weight-both},
\begin{multline*}
 w_{jk}-\log R\,\ell_{jk}=\frac{2\varepsilon^{(m)}_j\varepsilon^{(M)}_k}{mMR^2}
 \bigl\{R^2S_R+\log R[\cos(2R\theta_j)+\cos(2R\phi_k)]\\
 -2R^2h_m\cos(m\theta_j)\bigr\},
\end{multline*}
so, as $R^2|2h_m|<2\le2D_{jk}$,
\[
 |w_{jk}-\log R\,\ell_{jk}|\le\frac{112\,\varepsilon^{(m)}_j\varepsilon^{(M)}_kD_{jk}}{mMR^2}
 \le112\cdot\frac{mM}{R^2}\,\alpha_j\beta_k\le840\,\alpha_j\beta_k ,
\]
by Lemma~\ref{lem:tensor-cc} ($D_{jk}/(m^2M^2)\le\alpha_j\beta_k$). Multiplying by
$p(z_{jk})^2$, summing over the retained nodes, enlarging the sum to all
$(j,k)$ and using \eqref{eq:tensor-bound},
$|Q_m^\Pa(p^2)-\log R\,L_m[p]|\le840\cdot4\|p\|^2=3360\|p\|^2\le4096\|p\|^2$.
This is \eqref{eq:log-error} (the constants are not optimized).
Next, $|\ell_{jk}|\le4\varepsilon^{(m)}_j\varepsilon^{(M)}_k/(mMR^2)$, and every node has
$(1+d_j)(1+d'_k)\ge1$, so Lemma~\ref{lem:tail-estimate} with $K=1$ gives
\eqref{eq:L-bound}. Taking the infimum and the supremum of
\eqref{eq:log-error} over the unit sphere of $\mathcal P_{m-1}$, and using
$\log R>0$, gives the first inequality of \eqref{eq:log-spectral}, and
$|\nu_\pm|\le120$. Since $m/R\le5/2$,
\begin{align*}
 \Bigl|\frac{\lambda^\Pa_\pm}{\log m}-\nu_\pm\Bigr|
 &\le\frac{|\lambda^\Pa_\pm-\log R\,\nu_\pm|+|\nu_\pm|\log(m/R)}{\log m}\\
 &\le\frac{4096+120\log(5/2)}{\log m}<\frac{4300}{\log m}.\qquad\qed
\end{align*}

\subsection{Proof of Lemma~\ref{lem:corner-kernel}}\label{app:kernel}

\paragraph{Step 1: classical facts, with proofs}
Define $\mathsf P_k$ by Rodrigues' formula
$\mathsf P_k(x)=(2^kk!)^{-1}\frac{\mathrm d^k}{\mathrm dx^k}(x^2-1)^k$.
Expanding $(x^2-1)^k$ binomially and differentiating,
\[
 \mathsf P_k(x)=2^{-k}\sum_{r\le k/2}(-1)^r\frac{(2k-2r)!}{r!\,(k-r)!\,(k-2r)!}\,x^{k-2r}.
\]
This is the coefficient of $z^k$ in
$(1-z(2x-z))^{-1/2}=\sum_n\binom{2n}n4^{-n}z^n(2x-z)^n$ (put $r=k-n$). Thus
\begin{equation}\label{eq:legendre-gen}
 \sum_{k\ge0}\mathsf P_k(x)z^k=(1-2xz+z^2)^{-1/2}.
\end{equation}
Consequently $\mathsf P_k(1)=1$ and $\mathsf P_k(-x)=(-1)^k\mathsf P_k(x)$.
Integration by parts $k$ times shows that $\mathsf P_k$ is orthogonal to
polynomials of lower degree on $[-1,1]$, since the first $k-1$ derivatives of
$(x^2-1)^k$ vanish at $\pm1$. Its leading coefficient is
$(2k)!/[2^k(k!)^2]$, and
$\int_{-1}^1x^k\mathsf P_k\dd x=2^{-k}\int_{-1}^1(1-x^2)^k\dd x
=2^{-k}2^{2k+1}(k!)^2/(2k+1)!$ (integration by parts $k$ times, and the
recurrence for the Beta integral), so that
$\int_{-1}^1\mathsf P_k^2\dd x=2/(2k+1)$. Hence the functions
$\sqrt{(k+\frac12)(l+\frac12)}\,\mathsf P_k(x)\mathsf P_l(y)$, $k+l\le m-1$, form
an orthonormal basis of $\mathcal P_{m-1}$ in $L^2([-1,1]^2)$, and
\eqref{eq:repro-kernel} is its reproducing kernel.

For $x=\cos t$ and $\mathsf q=|\sin t|$ the geometric series gives, for small $|z|$,
$(1-zx-iz\mathsf q\cos\varphi)^{-1}=(1-zx)^{-1}\sum_n[iz\mathsf q\cos\varphi/(1-zx)]^n$.
Since $\frac1\pi\int_0^\pi\cos^{2r}\varphi\dd\varphi=\binom{2r}r4^{-r}$ and the odd
moments vanish,
\begin{align*}
 \sum_{k\ge0}z^k\frac1\pi\int_0^\pi(x+i\mathsf q\cos\varphi)^k\dd\varphi
 &=\frac1{1-zx}\Bigl(1+\frac{z^2\mathsf q^2}{(1-zx)^2}\Bigr)^{-1/2}\\
 &=(1-2xz+z^2)^{-1/2}.
\end{align*}
Comparison with \eqref{eq:legendre-gen}, and $\varphi\mapsto\pi-\varphi$ to
fix the sign of $\sin t$, give (Laplace's integral), for every real $t$,
\begin{equation}\label{eq:laplace}
 \mathsf P_k(\cos t)=\frac1\pi\int_0^\pi(\cos t+i\sin t\cos\varphi)^k\dd\varphi.
\end{equation}
Likewise, the exponential series and the same moments give
\begin{equation}\label{eq:bessel-rep}
 J_0(y)=\frac1\pi\int_0^\pi e^{iy\cos\varphi}\dd\varphi=\frac1\pi\int_0^\pi\cos(y\cos\varphi)\dd\varphi.
\end{equation}
Since $|\cos t+i\sin t\cos\varphi|\le1$, we get $|\mathsf P_k|\le1$ on $[-1,1]$,
and $|J_0|\le1$, $|J_0'|\le1$ on $\R$ (differentiate \eqref{eq:bessel-rep}).

\paragraph{Step 2: uniform approximation for $k\le m$}
Fix $c\in[-1,1]$ and let $X$ take the values $\pm1$ with probabilities
$(1\pm c)/2$, so $\mathbb EX=c$ and $\mathbb E(X-c)^2=1-c^2$. From
$|e^{iy}-1-iy|\le y^2/2$,
\[
 \bigl|\cos t+ic\sin t-e^{itc}\bigr|=\bigl|\mathbb E\,e^{itc}\bigl(e^{it(X-c)}-1-it(X-c)\bigr)\bigr|
 \le\tfrac{t^2}2(1-c^2).
\]
Both $\cos t+ic\sin t$ and $e^{itc}$ have modulus at most one, so the difference
of their $k$-th powers is at most $k$ times the difference. Taking
$c=\cos\varphi$, using \eqref{eq:laplace}, \eqref{eq:bessel-rep} (with $y=kt$) and
$\frac1\pi\int_0^\pi\sin^2\varphi\dd\varphi=\frac12$,
\begin{equation}\label{eq:legendre-bessel}
 |\mathsf P_k(\cos t)-J_0(kt)|\le kt^2/4\qquad(k\ge0,\ t\in\R),
\end{equation}
which at $k=0$ is the identity $0=0$. For $u\ge0$, $\varsigma\in[0,1]$ and
$0\le k\le m$, put $t=u/(m+\varsigma)$. Then \eqref{eq:legendre-bessel} and
$|J_0'|\le1$ give
\begin{equation}\label{eq:uniform-approx}
 \bigl|\mathsf P_k(\cos t)-J_0(ku/m)\bigr|\le\frac{ku^2}{4(m+\varsigma)^2}
 +\frac{ku\varsigma}{m(m+\varsigma)}\le\frac{u^2/4+\varsigma u}m .
\end{equation}

\paragraph{Step 3: reduction to a triangular sum}
By parity, $\mathsf P_k(s_1\cos t)=s_1^k\mathsf P_k(\cos t)$. Let $\xi_1=0$ if
$s_1=s_1'$ and $\xi_1=1$ otherwise, and $\xi_2=0$ if $s_2=s_2'$ and $\xi_2=1$
otherwise. With $w_{kl}=(k+\frac12)(l+\frac12)/m^2$,
\[
 m^{-4}\mathcal K_{m-1}(z_m,z_m')=m^{-2}\sum_{k+l\le m-1}w_{kl}(-1)^{\xi_1 k+\xi_2 l}\,
 \Pi_{kl},
\]
where $\Pi_{kl}$ is the product of the four factors
$\mathsf P_k(\cos\frac u{m+\varsigma_u})$, $\mathsf P_k(\cos\frac{u'}{m+\varsigma_{u'}})$,
$\mathsf P_l(\cos\frac v{m+\varsigma_v})$, $\mathsf P_l(\cos\frac{v'}{m+\varsigma_{v'}})$.
Let $F(a,b)=ab\,J_0(au)J_0(au')J_0(bv)J_0(bv')$. All factors involved have
modulus at most one, so \eqref{eq:uniform-approx} and telescoping give
\[
 \Bigl|\Pi_{kl}-J_0(\tfrac kmu)J_0(\tfrac kmu')J_0(\tfrac lmv)J_0(\tfrac lmv')\Bigr|
 \le\Theta/m .
\]
For $k+l\le m-1$ we have $(k+\frac12)+(l+\frac12)\le m$, hence
$0\le w_{kl}\le\frac14$; the number of terms is $m(m+1)/2\le m^2$; and
$|w_{kl}-\frac km\frac lm|=\frac{k+l}{2m^2}+\frac1{4m^2}\le\frac1{2m}$. Hence
\begin{equation}\label{eq:kernel-reduction}
 \Bigl|m^{-4}\mathcal K_{m-1}(z_m,z_m')-m^{-2}\sum_{k+l\le m-1}(-1)^{\xi_1 k+\xi_2 l}
 F(\tfrac km,\tfrac lm)\Bigr|\le\frac{\Theta/4+1/2}m .
\end{equation}

\paragraph{Step 4: equal corners}
On $[0,1]^2$, $|F|\le1$, $|\partial_aF|\le1+u+u'$ and $|\partial_bF|\le1+v+v'$
(use $|J_0|,|J_0'|\le1$). Let $\mathcal U_m$ be the union of the cells
$[\frac km,\frac{k+1}m]\times[\frac lm,\frac{l+1}m]$ with $k+l\le m-1$. It
contains $\mathcal T$ up to a null set, has area $(m+1)/(2m)\le1$, and
$\mathcal U_m\setminus\mathcal T$ has area $1/(2m)$. On each cell, $F$ differs from
$F(\frac km,\frac lm)$ by at most $(2+S)/m$; there are at most $m^2$ cells, each of
area $m^{-2}$, so
$|m^{-2}\sum_{k+l\le m-1}F(\frac km,\frac lm)-\int_{\mathcal U_m}F|\le(2+S)/m$.
Since $|F|\le1$, $|\int_{\mathcal U_m\setminus\mathcal T}F|\le1/(2m)$, whence
\begin{equation}\label{eq:riemann}
 \Bigl|m^{-2}\sum_{k+l\le m-1}F(\tfrac km,\tfrac lm)-\int_{\mathcal T}F\Bigr|\le\frac{5/2+S}m .
\end{equation}
For $\xi_1=\xi_2=0$, \eqref{eq:kernel-reduction} and \eqref{eq:riemann}
give \eqref{eq:kernel-limit}.

\paragraph{Step 5: different corners}
Let $\xi_1=1$. For fixed $l$ pair successive terms of
$\sum_{k=0}^{m-1-l}(-1)^kF(\frac km,\frac lm)$. By the derivative bound each
complete pair has modulus at most $(1+u+u')/m$; there are at most $m$ pairs, and
a possible unpaired last term has modulus at most one. So every row has modulus
at most $2+u+u'$, also after multiplication by the row sign $(-1)^{\xi_2 l}$.
Summing over the $m$ rows and dividing by $m^2$,
\begin{equation}\label{eq:alternating}
 \Bigl|m^{-2}\sum_{k+l\le m-1}(-1)^{k+\xi_2 l}F(\tfrac km,\tfrac lm)\Bigr|
 \le\frac{2+u+u'}m .
\end{equation}
If $\xi_1=0$ and $\xi_2=1$ use columns and obtain $(2+v+v')/m$. With
\eqref{eq:kernel-reduction}, the distance from $m^{-4}\mathcal K_{m-1}(z_m,z_m')$
to $0$ is at most $(\Theta/4+\frac52+S)/m\le(3+S+\Theta/4)/m$. This proves
\eqref{eq:kernel-limit}. The normalization is
$\int_{\mathcal T}ab\dd a\dd b=\frac12\int_0^1a(1-a)^2\dd a=\frac1{24}$. On the cube
$[0,U]^4$, $S\le4U$ and $\Theta\le U^2+4U$ (respectively $U^2+2U$ for the
Padua shifts), which gives the constants $C_U$. \qed

\subsection{Proof of Theorem~\ref{thm:corner-operator} and Remark~\ref{rem:finite-comparison}}\label{app:operator}

Fix $K\ge1$ and an integer $m\ge\max(4,2K)$. Let
$I_K=\{(r,s)\in\mathbb Z_{\ge0}^2:(1+r)(1+s)\le K\}$, so $r,s\le K-1$ and
$|I_K|\le K^2$.

\paragraph{Step 1: local indices, coefficients, active corners}
Since $K-1<m/2$, the nodes with $(1+d_j)(1+d'_k)\le K$ form four disjoint
corner patches. A \emph{patch node} is a symbol $\zeta=(\delta_x,\delta_y;r,s)$,
with $\delta_x,\delta_y\in\{0,1\}$ and $(r,s)\in I_K$, standing for the node with
$j=r$ if $\delta_x=0$ and $j=m-r$ if $\delta_x=1$, and with $k=s$ if
$\delta_y=0$ and $k=M-s$ if $\delta_y=1$. Then $d_j=r$, $d'_k=s$,
$\varepsilon^{(m)}_j=e_r$, $\varepsilon^{(M)}_k=e_s$, and
$z_{m,\zeta}=\bigl((-1)^{\delta_x}\cos\frac{\pi r}m,\,(-1)^{\delta_y}\cos\frac{\pi s}M\bigr)$.
The node is retained iff $r+s+m\delta_x+M\delta_y$ is even; this depends
only on the parity of $m$. Since $2R$ is even, the reflections
$j\mapsto m-j$, $k\mapsto M-k$ leave $\cos(2R\theta_j)$ and $\cos(2R\phi_k)$
unchanged, so the coefficient of $p(z_{m,\zeta})^2$ in \eqref{eq:L-m} is
$m^{-4}\gamma_m(\zeta)$, where
\[
 \gamma_m(\zeta)=-q_m\,e_re_s\Bigl[\cos\frac{2R\pi r}m+\cos\frac{2R\pi s}M\Bigr],\qquad
 q_m=\frac{2m^3}{MR^2}.
\]
As $m\to\infty$, $q_m\to8$ and the two cosines tend to $(-1)^r$ and $(-1)^s$.
For a retained node with $r+s$ even, $\gamma_m(\zeta)\to\gamma_\infty(\zeta):=-16e_re_s(-1)^r$;
for $r+s$ odd the two limits cancel and $\gamma_m(\zeta)\to\gamma_\infty(\zeta):=0$.
A retained node has $r+s\equiv m\delta_x+M\delta_y\pmod2$. So the corners
containing retained nodes with $r+s$ even (the \emph{active corners}) are those
with $m\delta_x+M\delta_y$ even: for even $m$ these are $\delta_y=0$, that is the
corners $(\pm1,1)$ of the square, and for odd $m$ they are $\delta_x=0$, the
corners $(1,\pm1)$. The corner $(1,1)$ is active for both parities. In the two
inactive corners every retained node has $r+s$ odd and limit coefficient zero.
Each active corner retains exactly the nodes with $(r,s)\in I_K$, $r+s$ even, and each
inactive corner those with $r+s$ odd, so the total number of retained patch nodes
is $N_K=2|I_K|$.

We also need quantitative control. For $m\ge4$: $q_m=8m/(m+1)$ for even $m$, and
$q_m-8=8(m^2+m-1)/[(m+1)(m-1)^2]\in(0,16/m]$ for odd $m\ge5$ (the upper bound
is equivalent to $m^3-3m^2-m+2\ge0$). Hence $|q_m-8|\le16/m$. The first cosine
differs from $(-1)^r$ by at most $\pi r/m$ and the second from $(-1)^s$ by at
most $2\pi s/m$ (as $M-2R\in\{1,2\}$). For an active node this gives
$|\gamma_m-\gamma_\infty|\le[16\cdot2+8\pi(r+2s)]/m\le(32+24\pi K)/m$. For an
inactive node, $|\gamma_m|\le q_m\cdot3\pi(K-1)/m\le(8+16/m)\,3\pi(K-1)/m
\le24\pi K/m$, using $m\ge2K$. Therefore the diagonal matrices
$C_m=\diag(\gamma_m(\zeta))$ and $C_\infty=\diag(\gamma_\infty(\zeta))$, on any
set of patch nodes, satisfy
\begin{equation}\label{eq:coeff-diff}
 \|C_m-C_\infty\|\le\beta_K(m):=\frac{32+24\pi K}m,\qquad\|C_\infty\|\le16 .
\end{equation}

\paragraph{Step 2: finite Gram kernels and finite-rank spectral transfer}
Let $\mathcal H_m=\mathcal P_{m-1}$ with the norm $\|\cdot\|$, and for a patch node
$\zeta$ let $v_{m,\zeta}=m^{-2}\mathcal K_{m-1}(\cdot,z_{m,\zeta})\in\mathcal H_m$.
By the reproducing property $p(z_{m,\zeta})=m^2\langle p,v_{m,\zeta}\rangle$, so
for a set $\mathcal Z$ of retained patch nodes
\[
 L^{\mathcal Z}_m[p]:=m^{-4}\sum_{\zeta\in\mathcal Z}\gamma_m(\zeta)p(z_{m,\zeta})^2
 =\sum_{\zeta\in\mathcal Z}\gamma_m(\zeta)\langle p,v_{m,\zeta}\rangle^2
\]
is the part of \eqref{eq:L-m} coming from $\mathcal Z$, and
$G_m(\zeta,\zeta')=\langle v_{m,\zeta},v_{m,\zeta'}\rangle=m^{-4}\mathcal K_{m-1}(z_{m,\zeta},z_{m,\zeta'})$.
Lemma~\ref{lem:corner-kernel} applies with $u=\pi r$, $v=\pi s$, $\varsigma_u=0$,
$\varsigma_v=1$ and the corner signs $((-1)^{\delta_x},(-1)^{\delta_y})$; as
$u,v\le U:=\pi(K-1)$, it gives
\begin{equation}\label{eq:gram-diff}
 |G_m(\zeta,\zeta')-G_\infty(\zeta,\zeta')|\le C_U/m,\qquad
 G_\infty(\zeta,\zeta')=\boldsymbol1_{(\delta_x,\delta_y)=(\delta_x',\delta_y')}\,
 \langle\psi_{rs},\psi_{r's'}\rangle ,
\end{equation}
because $\mathcal K_\infty((\pi r,\pi s),(\pi r',\pi s'))=\langle\psi_{rs},\psi_{r's'}\rangle$.

\begin{lemma}[Finite-rank transfer]\label{lem:finite-rank}
Let $m\ge\max(4,2K)$ have a fixed parity, let $\mathcal Z$ be a set of retained patch
nodes for that parity, $N=|\mathcal Z|\le N_K$, and, for a corner
$c=(\delta_x,\delta_y)$, $\mathcal Z_{c}=\{(r,s):(c;r,s)\in\mathcal Z\}$. Let
$B^{\mathcal Z}=\bigoplus_cB^{\mathcal Z_c}$ on $\mathcal H^{\oplus4}$, where
$B^{\mathcal Z_c}=\sum_{(r,s)\in\mathcal Z_c}\gamma_\infty(c;r,s)\langle\cdot,\psi_{rs}\rangle\psi_{rs}$
acts on $\mathcal H$. Let $\nu^{\mathcal Z}_-(m)$ and $\nu^{\mathcal Z}_+(m)$ be the
infimum and the supremum of $L^{\mathcal Z}_m[p]/\|p\|^2$ over $0\ne p\in\mathcal H_m$. Then
\[
 \bigl|\nu^{\mathcal Z}_-(m)-\inf\sigma(B^{\mathcal Z})\bigr|\le E_K(m),\qquad
 \bigl|\nu^{\mathcal Z}_+(m)-\sup\sigma(B^{\mathcal Z})\bigr|\le E_K(m),
\]
with $E_K$ as in Remark~\ref{rem:finite-comparison}. Moreover, the ordered
eigenvalues of the $N\times N$ matrix $X_m=G_m^{1/2}C_mG_m^{1/2}$ differ from
those of $X_\infty=G_\infty^{1/2}C_\infty G_\infty^{1/2}$ by at most $E_K(m)$,
the nonzero eigenvalues of $L_m^{\mathcal Z}$ are those of $X_m$, and the nonzero
eigenvalues of $B^{\mathcal Z}$ are those of $X_\infty$.
\end{lemma}
\begin{proof}
Let $V_m:\R^N\to\mathcal H_m$ send the unit vector of $\zeta$ to $v_{m,\zeta}$. Then
$L_m^{\mathcal Z}=V_mC_mV_m^*$ and $V_m^*V_m=G_m\ge0$. By the polar decomposition
$V_m=W_mG_m^{1/2}$ with a partial isometry $W_m$, the nonzero eigenvalues of
$V_mC_mV_m^*=W_m(G_m^{1/2}C_mG_m^{1/2})W_m^*$ are those of $X_m$. In the same way,
with $V_\infty:\R^N\to\mathcal H^{\oplus4}$ sending the unit vector of $\zeta$ to
$\psi_{rs}$ placed in the copy of $\mathcal H$ belonging to its corner, and with
\eqref{eq:gram-diff}, $B^{\mathcal Z}=V_\infty C_\infty V_\infty^*$ has as nonzero
eigenvalues those of $X_\infty$. Since $N\le N_K\le2K^2<K(2K+1)\le m(m+1)/2=\dim\mathcal H_m$,
zero is an eigenvalue of $L_m^{\mathcal Z}$, and $0\in\sigma(B^{\mathcal Z})$ because
$\dim\mathcal H=\infty$. Thus $\nu^{\mathcal Z}_-(m)=\min(0,\lambda_{\min}(X_m))$,
$\inf\sigma(B^{\mathcal Z})=\min(0,\lambda_{\min}(X_\infty))$, and similarly for the
suprema. It remains to bound $\|X_m-X_\infty\|$.

The entries of $G_m$ are bounded by the largest diagonal entry, which by
$|\mathsf P_k|\le1$ is at most the value of $m^{-4}\mathcal K_{m-1}$ at the corner
$(1,1)$:
$m^{-4}\sum_{k+l\le m-1}(k+\frac12)(l+\frac12)=m(m+1)(m^2+m+1)/(24m^4)<\frac18$
for $m\ge4$. (The closed form follows by induction on $m$: the increment from
$m-1$ to $m$ is $m(2m^2+1)/12$.) The entries of $G_\infty$ are at most
$\int_{\mathcal T}ab=\frac1{24}$. Hence $\|G_m\|,\|G_\infty\|\le N/8$, and by
\eqref{eq:gram-diff} $\delta:=\|G_m-G_\infty\|\le NC_U/m$.
For positive semidefinite matrices $A,B$ with $\|A-B\|=\delta$ we have
$\|A^{1/2}-B^{1/2}\|\le\frac4\pi\sqrt\delta\le2\sqrt\delta$. Indeed
$A^{1/2}=\pi^{-1}\int_0^\infty t^{-1/2}A(A+t)^{-1}\dd t$ (the scalar integral after
diagonalization). The difference of the integrands has norm at most $t^{-1/2}$,
since both $A(A+t)^{-1}$ and $B(B+t)^{-1}$ lie between $0$ and $I$, and, by the
resolvent identity $A(A+t)^{-1}-B(B+t)^{-1}=t(A+t)^{-1}(A-B)(B+t)^{-1}$, at most
$\delta t^{-3/2}$. Splitting the integral at $t=\delta$ gives the claim.
So $\|G_m^{1/2}-G_\infty^{1/2}\|\le2\sqrt{NC_U/m}$, and telescoping
$X_m-X_\infty$ with \eqref{eq:coeff-diff},
\begin{align*}
 \|X_m-X_\infty\|&\le\|G_m\|\,\|C_m-C_\infty\|
  +16\bigl(\|G_m^{1/2}\|+\|G_\infty^{1/2}\|\bigr)\|G_m^{1/2}-G_\infty^{1/2}\|\\
 &\le\frac{N\beta_K(m)}8+\frac{64}{\sqrt8}N\sqrt{C_U/m}\le E_K(m).
\end{align*}
By Weyl's inequality the ordered eigenvalues of $X_m$ and $X_\infty$ differ by at
most $\|X_m-X_\infty\|$, and $|\min(0,x)-\min(0,y)|\le|x-y|$.
\end{proof}

\paragraph{Step 3: the operator $B$ and the bound \eqref{eq:B-tail}}
Let $K'>K$. Apply Lemma~\ref{lem:finite-rank} with $K'$ in place of $K$ and
with $\mathcal Z=\{(0,0;r,s):K<(1+r)(1+s)\le K',\ r+s\text{ even}\}$ (nodes of the
corner $(1,1)$, active for both parities). Then $B^{\mathcal Z}=B_{K'}-B_K$ on the
first copy of $\mathcal H$ (the other copies vanish), because
$\gamma_\infty=-16e_re_s(-1)^r$ for these nodes. By the lemma, the extremes of
$L^{\mathcal Z}_m$ converge to those of $B_{K'}-B_K$ as $m\to\infty$, hence
$\|L^{\mathcal Z}_m\|\to\|B_{K'}-B_K\|$. All nodes of $\mathcal Z$ have
$(1+d_j)(1+d'_k)>K$ and $|\ell_{jk}|\le4\varepsilon^{(m)}_j\varepsilon^{(M)}_k/(mMR^2)$, so
Lemma~\ref{lem:tail-estimate} gives $\|L^{\mathcal Z}_m\|\le120/K$. Therefore
\begin{equation}\label{eq:B-cauchy}
 \|B_{K'}-B_K\|\le120/K\qquad(K'>K).
\end{equation}
So $(B_K)$ is a Cauchy sequence of selfadjoint operators and converges in norm to a
selfadjoint operator $B$. It is compact, being a norm limit of finite-rank operators,
and letting $K'\to\infty$ in \eqref{eq:B-cauchy} gives \eqref{eq:B-tail}.

\paragraph{Step 4: the limits \eqref{eq:corner-limits}}
Put $c_-(K)=\inf\sigma(B_K)$ and $c_+(K)=\sup\sigma(B_K)$ (both include $0$);
by \eqref{eq:B-tail}, $|c_\pm(K)-c_\pm|\le120/K$. Let $\mathcal Z_K$ be the set of
all retained patch nodes for the parity of $m$ and $L_m^K=L_m^{\mathcal Z_K}$, the part of
\eqref{eq:L-m} from the nodes with $(1+d_j)(1+d'_k)\le K$. By Step 1, the two
active corners carry the same retained index set
$\{(r,s)\in I_K:r+s\text{ even}\}$ and the same coefficients $\gamma_\infty$, while
the two inactive corners have $\gamma_\infty=0$. Hence
$B^{\mathcal Z_K}=B_K\oplus B_K\oplus0\oplus0$ (up to the order of the corners),
whose extremes are those of $B_K$. Let $\nu^K_\pm(m)$ be the extremes of
$L_m^K$. Lemma~\ref{lem:finite-rank} gives
$\nu^K_\pm(m)\to c_\pm(K)$ as $m\to\infty$ along each parity, and
Lemma~\ref{lem:tail-estimate} gives
$|\nu_\pm(m)-\nu^K_\pm(m)|\le\|L_m-L_m^K\|\le120/K$. Consequently
$\limsup_{m\to\infty}|\nu_\pm(m)-c_\pm|\le240/K$ along each parity, and
$K\to\infty$ shows $\nu_\pm(m)\to c_\pm$ over all integers. By
\eqref{eq:log-spectral}, $\lambda^\Pa_\pm(m,m-1)/\log m\to c_\pm$.

Finally, Proposition~\ref{prop:carrier-improvement} gives
$\limsup\lambda^\Pa_-/\log m\le-\kappa_{\mathrm{new}}$, hence
$c_-\le-\kappa_{\mathrm{new}}<0$, and \eqref{eq:padua-growth-bounds} gives
$c_+\ge\kappa^{\mathrm{car}}_+>0$. These two inequalities serve only for the
signs; the identification of the limits above does not use any carrier. \qed

\paragraph{Step 5: proof of \eqref{eq:finite-comparison}}
For $m\ge\max(4,2K)$, Lemma~\ref{lem:finite-rank} with $\mathcal Z_K$ gives
$|\nu^K_\pm(m)-c_\pm(K)|\le E_K(m)$. Combining this with
\eqref{eq:log-spectral}, $|\nu_\pm-\nu^K_\pm|\le120/K$ and
$|c_\pm(K)-c_\pm|\le120/K$,
\[
 \Bigl|\frac{\lambda^\Pa_\pm}{\log m}-c_\pm\Bigr|
 \le\frac{4300}{\log m}+\frac{120}K+E_K(m)+\frac{120}K.\qquad\qed
\]

\paragraph{Proof of Remark~\ref{rem:two-corners}}
It follows from Steps 1 and 4: $X_\infty=X_B\oplus X_B\oplus0\oplus0$ with the
same $X_B$ for the two active corners, whose nonzero eigenvalues are those of
$B_K$; Lemma~\ref{lem:finite-rank} gives the convergence, with multiplicity, of
the eigenvalues of $X_m$ to those of $X_\infty$.

\section{Proofs for the structure and enclosure of \texorpdfstring{$B$}{B}}\label{app:structure}

Throughout, $\mathcal T$, $\mathcal H$, $\dd\mu=ab\dd a\dd b$, $B_K$, $B$, $\psi_{jk}$, $e_j$ and
$c_\pm$ are as in Theorem~\ref{thm:corner-operator} and Section~\ref{sec:structure},
$B^-=-B$ and $r=\|B\|$. All functions are real valued. We use the integral
representation \eqref{eq:bessel-rep} of $J_0$ and $|J_0|\le1$
(Appendix~\ref{app:kernel}). Unless stated otherwise, integrals over $(0,1)$
carry the weight $a\dd a$.

\subsection{Arcsine densities and the scalar kernels}\label{app:arcsine}

For $a>0$ let $\omega_a(s)=\boldsymbol1_{|s|<a}/(\pi\sqrt{a^2-s^2})$. It is
a probability density on $\R$, namely the law of $a\cos\alpha$ for $\alpha$
uniform on $[0,2\pi)$. Substituting $s=a\cos\theta$ and using \eqref{eq:bessel-rep},
\begin{equation}\label{eq:omega-ft}
 \int_\R e^{ixs}\omega_a(s)\dd s=\frac1\pi\int_0^\pi e^{iax\cos\theta}\dd\theta=J_0(ax).
\end{equation}
For $0<a,a'<1$ let $q=q_{a,a'}=\omega_a*\omega_{a'}$, the density of
$a\cos\alpha+a'\cos\alpha'$ for independent uniform $\alpha,\alpha'$. It is
nonnegative, even, supported in $[-(a+a'),a+a']\subset(-2,2)$, and by
\eqref{eq:omega-ft} its Fourier transform is $J_0(ax)J_0(a'x)$. Let
$q^\circ(s)=\int\omega_a(u)\omega_{a'}(s-u)\dd u\in[0,\infty]$ be the convolution integral
itself, defined for every $s$, which is a version of $q$. Put
$\xi_s^2=((a+a')^2-s^2)/(4aa')$.

\begin{lemma}[Arcsine convolution]\label{lem:arcsine}
Let $0<a,a'<1$ with $a\ne a'$.
\begin{enumerate}
\item[(i)] $q^\circ$ is finite and continuous on each of the open sets
 $\{|s|<|a-a'|\}$, $\{|a-a'|<|s|<a+a'\}$, $\{|s|>a+a'\}$, and on them
 \[
  q^\circ(s)=\begin{cases}
  \dfrac{\mathbf K(1/\xi_s)}{\pi^2\sqrt{aa'}\,\xi_s}&|s|<|a-a'|,\\[2ex]
  \dfrac{\mathbf K(\xi_s)}{\pi^2\sqrt{aa'}}&|a-a'|<|s|<a+a',\\[1ex]
  0&|s|>a+a'.\end{cases}
 \]
\item[(ii)] With $\Sigma_0,\Sigma_1$ of \eqref{eq:sigma-kernels},
 $q^\circ(0)=\Sigma_0(a,a')$, and $2q^\circ(1)=\Sigma_1(a,a')$ if $a+a'\ne1$.
\end{enumerate}
\end{lemma}
\begin{proof}
Let $\phi=\alpha'-\alpha$ modulo $2\pi$ and write
$a+a'e^{i\phi}=\varpi(\phi)e^{i\gamma(\phi)}$ with $\varpi(\phi)^2=a^2+a'^2+2aa'\cos\phi$.
Then $a\cos\alpha+a'\cos\alpha'=\operatorname{Re}\bigl(e^{i\alpha}(a+a'e^{i\phi})\bigr)
=\varpi(\phi)\cos(\alpha+\gamma(\phi))$. The pair $(\alpha,\phi)$ is uniform on the
torus, so given $\phi$ the angle $\alpha+\gamma(\phi)$ is uniform and the conditional
density of the sum is $\omega_{\varpi(\phi)}$. Hence $q=\frac1{2\pi}\int_0^{2\pi}\omega_{\varpi(\phi)}\dd\phi$
almost everywhere. Folding $\phi\mapsto2\pi-\phi$ and then substituting
$\phi=\pi-\varphi$ gives, for almost every $s$,
\[
 q(s)=\frac1\pi\int_0^\pi\omega_{\varpi_-(\varphi)}(s)\dd\varphi
 =\frac1{\pi^2}\int_{\{\varpi_-(\varphi)>|s|\}}\frac{\dd\varphi}{\sqrt{\varpi_-(\varphi)^2-s^2}},
\]
where $\varpi_-(\varphi)^2=a^2+a'^2-2aa'\cos\varphi$.
Now substitute $\varphi=\pi-2u$ with $u\in(0,\pi/2)$. Then
$\varpi_-^2=(a+a')^2-4aa'\sin^2u$, $\varpi_-^2-s^2=4aa'(\xi_s^2-\sin^2u)$ and $\dd\varphi=2\dd u$, so
\[
 q(s)=\frac1{\pi^2\sqrt{aa'}}\int_{\{u\in(0,\pi/2):\ \sin u<\xi_s\}}\frac{\dd u}{\sqrt{\xi_s^2-\sin^2u}} .
\]
If $|a-a'|<|s|<a+a'$ then $0<\xi_s<1$, and the substitution $\sin u=\xi_s\sin v$
turns the integral into $\mathbf K(\xi_s)$. If $|s|<|a-a'|$ then $\xi_s>1$, every $u$
qualifies, and $\sin u=\sin v$ turns the integral into $\xi_s^{-1}\mathbf K(1/\xi_s)$. If
$|s|>a+a'$ the domain is empty. The three expressions are continuous in $s$ on the
three open sets. The function $q^\circ$ is also finite and continuous there: its
integrand has inverse square-root singularities only at $u=\pm a$ and
$u=s\pm a'$, and for $s$ in a compact subset of one of the three sets these four points
are pairwise separated (they can coincide only if $|s|\in\{|a-a'|,a+a'\}$);
splitting $\R$ into small neighbourhoods of these points, in each of which only one
factor is singular and the other is bounded and continuous, and a remainder on which
both factors are bounded, dominated convergence gives the continuity. Two continuous
functions that agree almost everywhere on an open set agree everywhere on it;
this proves (i).

(ii) For $s=0$ the first case applies, with $\xi_0=(a+a')/(2\sqrt{aa'})$, so that
$\sqrt{aa'}\,\xi_0=(a+a')/2$ and $1/\xi_0=k_0$; this gives $q^\circ(0)=\Sigma_0(a,a')$.
For $s=1$: if $a+a'<1$ the third case gives $q^\circ(1)=0=\Sigma_1/2$; if $a+a'>1$,
then $|a-a'|<1<a+a'$ and the second case applies with $\xi_1=k_1$, giving
$q^\circ(1)=\mathbf K(k_1)/(\pi^2\sqrt{aa'})=\Sigma_1/2$.
\end{proof}

\paragraph{Poisson regularization}
For $0<\varrho<1$ let $P_\varrho(s)=\sum_{n\in\mathbb Z}\varrho^{|n|}e^{i\pi ns}
=(1-\varrho^2)/(1-2\varrho\cos\pi s+\varrho^2)>0$. It is $2$-periodic, satisfies
$\frac12\int_{-1}^1P_\varrho=1$, and $\sup_{\delta\le|s|\le1}P_\varrho(s)\to0$
as $\varrho\uparrow1$ for each $\delta\in(0,1)$. For a $2$-periodic $F\in L^1_{\mathrm{loc}}$
set $(P_\varrho*F)(s)=\frac12\int_{-1}^1P_\varrho(s-u)F(u)\dd u
=\sum_n\varrho^{|n|}\widehat F(n)e^{i\pi ns}$, with
$\widehat F(n)=\frac12\int_{-1}^1F(u)e^{-i\pi nu}\dd u$. If $F$ coincides almost
everywhere near $s$ with a function that is continuous at $s$, then
$(P_\varrho*F)(s)\to F(s)$ as $\varrho\uparrow1$ (approximate identity: the
contribution of $\{|u-s|\ge\delta\}$ tends to zero, and that of the rest is
controlled by the continuity at $s$). For $t\in\{0,1\}$ put
\[
 \Sigma_t^\varrho(a,a')=\sum_{n\ge0}e_n(-1)^{tn}\varrho^nJ_0(\pi na)J_0(\pi na'),
\]
an absolutely convergent series. Let $q^\sharp=\sum_{l\in\mathbb Z}q(\cdot+2l)$ be the
$2$-periodization of $q=q_{a,a'}$. Since $\widehat{q^\sharp}(n)=\frac12\int_\R q(s)e^{-i\pi ns}\dd s
=\frac12J_0(\pi na)J_0(\pi na')$, and since the terms $\pm n$ combine,
\begin{equation}\label{eq:abel-identity}
 \Sigma_t^\varrho(a,a')=(P_\varrho*q^\sharp)(t)\ \ge0 .
\end{equation}

\begin{lemma}[Abel sums]\label{lem:abel}
For $t\in\{0,1\}$: (i) $\Sigma^\varrho_t(a,a')\to\Sigma_t(a,a')$ as $\varrho\uparrow1$ for
almost every $(a,a')\in(0,1)^2$; (ii) $\iint\Sigma^\varrho_t\,a\,a'\dd a\dd a'\to\iint\Sigma_t\,a\,a'\dd a\dd a'<\infty$;
(iii) $\Sigma_t^\varrho\to\Sigma_t$ in $L^1((0,1)^2,a\,a'\dd a\dd a')$.
\end{lemma}
\begin{proof}
(i) Let $a\ne a'$, and $a+a'\ne1$ if $t=1$. Near $s=0$ only the term $l=0$ of $q^\sharp$ is
nonzero, and near $s=1$ only $q(s)+q(s-2)$. By Lemma~\ref{lem:arcsine}(i), $q$ is
continuous near $0$ and near $\pm1$ (note $|a-a'|<1$, and $a+a'\ne1$). Hence by
\eqref{eq:abel-identity}, $\Sigma^\varrho_0\to q(0)$ and $\Sigma^\varrho_1\to q(1)+q(-1)=2q(1)$,
which are $\Sigma_0$ and $\Sigma_1$ by Lemma~\ref{lem:arcsine}(ii).

(ii) Let $\Omega(s)=\int_0^1\omega_a(s)\,a\dd a=\sqrt{(1-s^2)_+}/\pi$ (integrate
$a\dd a/(\pi\sqrt{a^2-s^2})$ over $a>|s|$). It is bounded by $1/\pi$ and integrable, so
$Z=\Omega*\Omega$ is continuous, even and supported in $[-2,2]$, with $Z(\pm2)=0$. By Tonelli,
$\iint q^\circ_{a,a'}(s)\,a\,a'\dd a\dd a'=Z(s)$ for every $s$. Hence
$\iint q^\sharp_{a,a'}(s)\,a\,a'\dd a\dd a'=Z^\sharp(s):=\sum_lZ(s+2l)$, a continuous function with
$Z^\sharp(0)=Z(0)$ and $Z^\sharp(1)=2Z(1)$. By \eqref{eq:abel-identity} and Tonelli,
\[
 \iint\Sigma^\varrho_t\,a\,a'\dd a\dd a'=(P_\varrho*Z^\sharp)(t)\longrightarrow Z^\sharp(t)\qquad(\varrho\uparrow1).
\]
On the other hand $\Sigma_0=q^\circ(0)$ and $\Sigma_1=2q^\circ(1)$ almost everywhere, so by
Tonelli $\iint\Sigma_0\,a\,a'=Z(0)$ and $\iint\Sigma_1\,a\,a'=2Z(1)$, the same limits.

(iii) Nonnegative functions that converge almost everywhere with converging finite integrals
converge in $L^1$ (Scheffé's lemma: $\int|u_\varrho-u|=\int u_\varrho+\int u-2\int\min(u_\varrho,u)$, and
the last integral tends to $\int u$ by dominated convergence). Apply this along any sequence
$\varrho_k\uparrow1$, using (i), (ii) and $\Sigma_t^\varrho\ge0$.
\end{proof}

\subsection{Proof of Proposition~\ref{prop:elliptic-kernel}}\label{app:elliptic-proof}

\paragraph{Step 1: measurability and row bounds}
Off the null sets $a=a'$ and $a+a'=1$ the functions $\Sigma_0,\Sigma_1$ are given by
explicit continuous expressions, so they are measurable, nonnegative ($\mathbf K>0$) and
symmetric. By Lemma~\ref{lem:arcsine}(ii) and Tonelli, for almost every $a$,
\begin{align*}
 \int_0^1\Sigma_0(a,a')\,a'\dd a'&=\int\omega_a(s)\Omega(-s)\dd s\le\frac1\pi,\\
 \int_0^1\Sigma_1(a,a')\,a'\dd a'&=2\int\omega_a(s)\Omega(1-s)\dd s\le\frac2\pi,
\end{align*}
because $\omega_a$ has mass one and $\Omega\le1/\pi$. This is \eqref{eq:row-bounds}.

\paragraph{Step 2: a bounded operator with the kernel}
Put $\tilde\beta(x,y)=8[\Sigma_1(a,a')\Sigma_0(b,b')+\Sigma_0(a,a')\Sigma_1(b,b')]\ge0$ for
$x=(a,b)$, $y=(a',b')\in(0,1)^2$. It is symmetric and, by \eqref{eq:row-bounds},
$\int_{(0,1)^2}\tilde\beta(x,y)\dd\mu(y)\le8[\frac2\pi\cdot\frac1\pi+\frac1\pi\cdot\frac2\pi]=32/\pi^2$
for almost every $x$. For $f,g\in L^2((0,1)^2,\dd\mu)$ the inequality
$|g(x)f(y)|\le\frac12(g(x)^2+f(y)^2)$ and the symmetry give
$\iint\tilde\beta|g(x)||f(y)|\dd\mu\dd\mu\le(32/\pi^2)\|f\|\|g\|$ (Schur's test). Hence $\tilde\beta$
defines a bounded selfadjoint integral operator $\tilde B$ on $\mathcal H$ (extend by zero and
compress to $\mathcal T$), with kernel $-\tilde\beta=\beta$ and $\|\tilde B\|\le32/\pi^2$.

\paragraph{Step 3: $B=\tilde B$}
Let $f,g\in C_c^\infty(\mathcal T^\circ)$, where $\mathcal T^\circ$ is the open triangle, and
$c_{jk}(h)=\int h\psi_{jk}\dd\mu$. The operator $\mathcal L_a=a^{-1}\partial_a(a\partial_a)$
satisfies $\mathcal L_aJ_0(\pi ja)=-(\pi j)^2J_0(\pi ja)$ (termwise differentiation of the
power series) and is symmetric for the weight $a\dd a$ on functions with compact support in
$(0,1)$. Integrating by parts $\rho_1$ times in $a$ and $\rho_2$ times in $b$,
\begin{equation}\label{eq:coeff-decay}
 |c_{jk}(h)|\le C_{h,\rho_1,\rho_2}(1+j)^{-2\rho_1}(1+k)^{-2\rho_2}\qquad(\rho_1,\rho_2\in\mathbb N),
\end{equation}
so $\sum_{j,k}|c_{jk}(f)c_{jk}(g)|<\infty$. Since $B_K\to B$ in norm,
\begin{equation}\label{eq:B-sum}
 \langle g,Bf\rangle=\lim_K\langle g,B_Kf\rangle
 =-16\sum_{j+k\text{ even}}e_je_k(-1)^jc_{jk}(f)c_{jk}(g),
\end{equation}
an absolutely convergent sum. For $0<\varrho<1$ the damped sum, with $\varrho^{j+k}$ inserted,
tends to the same limit as $\varrho\uparrow1$, by absolute convergence. Since
$\boldsymbol1_{j+k\text{ even}}(-1)^j=\frac12[(-1)^j+(-1)^k]$ and
$c_{jk}(f)c_{jk}(g)=\iint J_0(\pi ja)J_0(\pi ja')J_0(\pi kb)J_0(\pi kb')g(x)f(y)\dd\mu(x)\dd\mu(y)$,
the damped sum equals, after exchanging the (absolutely convergent) sums with the integral,
\[
 -8\iint\bigl\{\Sigma^\varrho_1(a,a')\Sigma^\varrho_0(b,b')+\Sigma^\varrho_0(a,a')\Sigma^\varrho_1(b,b')\bigr\}
 g f\dd\mu\dd\mu,
\]
where $g=g(x)$, $f=f(y)$, $x=(a,b)$ and $y=(a',b')$.
By Lemma~\ref{lem:abel}(iii) and Fubini, the products of $L^1$-convergent factors in the
independent variable pairs $(a,a')$, $(b,b')$ converge in $L^1(\dd\mu\otimes\dd\mu)$
($\|u_\varrho v_\varrho-uv\|_1\le\|u_\varrho-u\|_1\|v_\varrho\|_1+\|u\|_1\|v_\varrho-v\|_1$ with bounded
norms). As $g(x)f(y)$ is bounded, the displayed integral tends to
$\iint\beta(x,y)g(x)f(y)\dd\mu\dd\mu=\langle g,\tilde Bf\rangle$. Hence $\langle g,Bf\rangle=\langle g,\tilde Bf\rangle$ for
all such $f,g$. Both operators are bounded and $C_c^\infty(\mathcal T^\circ)$ is dense in
$\mathcal H$, so $B=\tilde B$. \qed

\subsection{Proof of Proposition~\ref{prop:positivity}}\label{app:positivity-proof}

The operator $B^-$ is compact and selfadjoint (Theorem~\ref{thm:corner-operator}) with
kernel $k^-=-\beta=\tilde\beta\ge0$ (Proposition~\ref{prop:elliptic-kernel}).
By Lemma~\ref{lem:arcsine}, $\Sigma_0>0$ off the diagonal and $\Sigma_1>0$ for $a+a'>1$. So, for
almost every $(x,y)$, $k^-(x,y)>0$ whenever $a+a'>1$ or $b+b'>1$. By the row bounds, for
$f\in\mathcal H$ the integral $\int k^-(x,y)|f(y)|\dd\mu(y)$ is finite for almost every $x$.

\paragraph{Proof of (i)}
Let $f\ge0$, $f\ne0$. For some $\delta>0$ the set $S=\{y\in\mathcal T:a'>\delta,\ f(y)>0\}$ has positive
measure. If $z=(s,t)\in\mathcal T$ and $s>1-\delta$, then $s+a'>1$ for all $y\in S$, so
$k^-(z,y)>0$ for almost every $y\in S$ and $B^-f(z)\ge\int_Sk^-(z,y)f(y)\dd\mu(y)>0$.
Fix $x=(a,b)$ with $a>0$ and let $W_x=\{(s,t)\in\mathcal T:s>\max(1-\delta,1-a)\}$, a set of
positive measure. For $z\in W_x$ we have $a+s>1$, hence $k^-(x,z)>0$ almost everywhere on $W_x$
and $B^-f>0$ on $W_x$. Since $k^-\ge0$ and $B^-f\ge0$,
$(B^-)^2f(x)\ge\int_{W_x}k^-(x,z)B^-f(z)\dd\mu(z)>0$. This holds for almost every $x$. In
particular $B^-\ne0$.

\paragraph{Proof of (ii)}
As $B^-\ne0$ is compact and selfadjoint, there is a real $v\ne0$ with $B^-v=\lambda v$ and
$|\lambda|=r>0$. Since $k^-\ge0$,
$r\|v\|^2=|\langle v,B^-v\rangle|\le\langle|v|,B^-|v|\rangle\le r\|v\|^2$. Thus
$u=|v|$ satisfies $\langle u,(r-B^-)u\rangle=0$. The operator $r-B^-$ is positive semidefinite,
so $(r-B^-)u=0$, i.e.\ $B^-u=ru$ with $u\ge0$, $u\ne0$. Then $(B^-)^2u=r^2u$ and (i)
gives $u>0$ almost everywhere. So $r=\sup\sigma(B^-)$, that is $c_-=\inf\sigma(B)=-r$.

\paragraph{Proof of (iii)}
Let $T=(B^-)^2$ and let $E$ be its real eigenspace for $r^2=\|T\|$; it contains $u>0$.
Suppose $\dim E\ge2$. Then there is $w\in E$, $w\ne0$, $\langle w,u\rangle=0$. As $u>0$ almost
everywhere, $w_\pm=\max(\pm w,0)$ are both nonzero. By (i), $Tw_->0$ almost everywhere and
$\langle w_+,Tw_-\rangle>0$. Now
$\langle|w|,T|w|\rangle-\langle w,Tw\rangle=4\langle w_+,Tw_-\rangle>0$, whereas
$\langle w,Tw\rangle=r^2\|w\|^2\ge\langle|w|,T|w|\rangle$ because $\|T\|=r^2$. This is a
contradiction, so $\dim E=1$. Eigenvectors of $B^-$ for $r$ and, if it were an eigenvalue,
for $-r$ are orthogonal elements of $E$; hence $r$ is a simple eigenvalue of $B^-$ and $-r$ is
not an eigenvalue of $B^-$. For $B=-B^-$: $-r$ is a simple eigenvalue with eigenfunction $u>0$,
unique up to a scalar factor, and $r$ is not an eigenvalue of $B$. By
Theorem~\ref{thm:corner-operator}, $c_+>0$ and $c_+\le\|B\|=r$; if $c_+=r$, then $r$,
a nonzero point of the spectrum of the compact operator $B$, would be an
eigenvalue, which was just excluded. Hence $0<c_+<r$. \qed

\subsection{Proof of Proposition~\ref{prop:rectangles}}\label{app:rect-proof}

For $0\le u<v\le1$ let $E(s)=\bigl[\sqrt{(v^2-s^2)_+}-\sqrt{(u^2-s^2)_+}\bigr]/\pi$. It is
continuous and even. For $|s|<v$, $\int_u^v\omega_{a'}(s)\,a'\dd a'=\frac1\pi\int_{\max(u,|s|)}^va'(a'^2-s^2)^{-1/2}\dd a'=E(s)$,
and the left side is $0=E(s)$ for $|s|\ge v$. On $|s|\le u$, $E$ is nondecreasing in $|s|$
(the derivative of $\sqrt{v^2-\sigma}-\sqrt{u^2-\sigma}$ in $\sigma=s^2$ is positive since
$u<v$) with value $\sqrt{v^2-u^2}/\pi$ at $|s|=u$; on $u\le|s|\le v$ it equals
$\sqrt{v^2-s^2}/\pi$, which decreases. Hence $0\le E\le\sqrt{v^2-u^2}/\pi$.

\paragraph{(a)}
By the substitution $s=a\cos\theta$,
$\chi_t\int\omega_a(s)E(t-s)\dd s=\frac{\chi_t}\pi\int_0^\pi E(t-a\cos\theta)\dd\theta=I_t(a;[u,v])$,
which is continuous in $a\in[0,1]$ (dominated convergence) and, since $\omega_a$ has mass one,
lies in $[0,\chi_t\sqrt{v^2-u^2}/\pi]$. For the identity, fix $a\in(0,1)$ and let
$F_a=\omega_a*E$, a continuous even function supported in $[-2,2]$ with
$F_a(\pm2)=0$. By Tonelli, for almost every $s$,
$\int_u^vq^\sharp_{a,a'}(s)\,a'\dd a'=F_a^\sharp(s):=\sum_lF_a(s+2l)$, which is continuous with
$F^\sharp_a(0)=F_a(0)$ and $F^\sharp_a(1)=2F_a(1)$. By \eqref{eq:abel-identity},
\[
 \int_u^v\Sigma_t^\varrho(a,a')\,a'\dd a'=(P_\varrho*F_a^\sharp)(t)\longrightarrow\chi_tF_a(t)=I_t(a;[u,v])
 \qquad(\varrho\uparrow1)
\]
for every $a$. By Lemma~\ref{lem:abel}(iii) the functions
$a\mapsto\int_u^v\Sigma_t^\varrho(a,a')\,a'\dd a'$ converge in $L^1((0,1),a\dd a)$ to
$a\mapsto\int_u^v\Sigma_t(a,a')\,a'\dd a'$, hence almost everywhere along a subsequence.
Therefore $I_t(a;[u,v])=\int_u^v\Sigma_t(a,a')\,a'\dd a'$ for almost every $a$. For a
rectangle $U\times V$, Tonelli and the product form of $-\beta$ give
$\int_{U\times V}(-\beta)(x,y)\dd\mu(y)=8\bigl[\int_U\Sigma_1(a,a')a'\dd a'\int_V\Sigma_0(b,b')b'\dd b'
+\int_U\Sigma_0(a,a')a'\dd a'\int_V\Sigma_1(b,b')b'\dd b'\bigr]=P_{U\times V}(x)$ for almost every $x=(a,b)$.

\paragraph{(b) Upper certificate}
Let $h=\sum_ih_i\boldsymbol1_{C_i}>0$ almost everywhere on $\mathcal T$ (the
rectangles not meeting $\mathcal T$ in a set of positive measure play no role). Since
$-\beta\ge0$ and $C_j\subset\mathcal R_j$, (a) gives for almost every $x\in C_i$
\[
 (B^-h)(x)=\sum_jh_j\int_{C_j}(-\beta)(x,y)\dd\mu(y)\le\sum_jh_jP_{\mathcal R_j}(x)\le Q_i\le\Lambda\,h_i,
\]
where $\Lambda=\max_lQ_l/h_l$.
For $f\in\mathcal H$, by $|f(x)||f(y)|\le\frac12\bigl[f(x)^2h(y)/h(x)+f(y)^2h(x)/h(y)\bigr]$ and the
symmetry of $-\beta$,
\begin{multline*}
 |\langle f,B^-f\rangle|\le\iint(-\beta)(x,y)|f(x)||f(y)|\dd\mu\dd\mu\\
 \le\int f(x)^2\frac{(B^-h)(x)}{h(x)}\dd\mu(x)\le\Lambda\|f\|^2 .
\end{multline*}
As $B^-$ is selfadjoint, $\|B\|=\sup_{\|f\|=1}|\langle f,B^-f\rangle|\le\Lambda$.

\paragraph{(c) Lower certificate}
With $h$ as above, $\|B\|\ge\langle h,B^-h\rangle/\|h\|^2$ and $\|h\|^2=\sum_ih_i^2\mu(C_i)$. Since
$-\beta\ge0$, omitting the terms with $i\notin F$ or $j\notin F$ decreases the form, and
$C_i=\mathcal R_i$ for $i\in F$, so
$\langle h,B^-h\rangle\ge\sum_{i,j\in F}h_ih_j\langle\boldsymbol1_{\mathcal R_i},B^-\boldsymbol1_{\mathcal R_j}\rangle$.
For $i,j\in F$, by (a) and Fubini,
\begin{multline*}
 \langle\boldsymbol1_{\mathcal R_i},B^-\boldsymbol1_{\mathcal R_j}\rangle=\int_{\mathcal R_i}P_{\mathcal R_j}\dd\mu\\
 =8\bigl[J_1(U_i,U_j)J_0(V_i,V_j)+J_0(U_i,U_j)J_1(V_i,V_j)\bigr]\ge\mathsf m_{ij}. \qed
\end{multline*}

\subsection{The finite check behind Proposition~\ref{prop:enclosure}}\label{app:certificate}

Take $\mathsf N=16$ and $\mathsf T=256$, the bins $W_i=[i/\mathsf N,(i+1)/\mathsf N]$ ($0\le i<\mathsf N$),
the squares $\mathcal R_{ij}=W_i\times W_j$ and $h\equiv1$ in
Proposition~\ref{prop:rectangles}. The squares meeting $\mathcal T$ in a set of positive
measure (outer cells) are those with $i+j\le\mathsf N-1$, and those contained in
$\mathcal T$ (inner cells) have $i+j\le\mathsf N-2$; there are $136$ and $120$
of them. Moreover $\|h\|^2=\mu(\mathcal T)=\frac12\int_0^1a(1-a)^2\dd a=\frac1{24}$ and
$\int_{W_i}a\dd a=(2i+1)/(2\mathsf N^2)$.

\paragraph{Enclosures of $I_t$}
Let $0\le u<v\le1$, $a\in W_i$, and let $g(\theta)=\sqrt{(v^2-z^2)_+}-\sqrt{(u^2-z^2)_+}$ with
$z=t-a\cos\theta$, so that $I_t(a;[u,v])=\frac{\chi_t}{\pi^2}\int_0^\pi g\dd\theta$ by
\eqref{eq:I-t}. By the proof of Proposition~\ref{prop:rectangles}(a),
$0\le g\le\sqrt{v^2-u^2}$. Split $[0,\pi]$ into $\mathsf T$ cells
$\Theta_l=[l\pi/\mathsf T,(l+1)\pi/\mathsf T]$. As $\cos$ decreases on $[0,\pi]$, for $\theta\in\Theta_l$
the value $\cos\theta$ lies in the interval $c_l$ between $\cos((l+1)\pi/\mathsf T)$ and
$\cos(l\pi/\mathsf T)$, so $z$ lies in $t-W_i\cdot c_l$ and $z^2$ in the exact hull $Z_l$ of the squares of
that interval (which contains $0$ if the interval does). The functions $\sigma\mapsto\sqrt{(v^2-\sigma)_+}$
and $\sigma\mapsto\sqrt{(u^2-\sigma)_+}$ are monotone, so $g(\theta)$ lies in the interval $G_l$
obtained by subtracting the interval images of $Z_l$ and intersecting with $[0,\sqrt{v^2-u^2}]$. Since
$\int_0^\pi g\dd\theta\in\frac\pi{\mathsf T}\sum_lG_l$,
\begin{equation}\label{eq:I-enclosure}
 I_t(a;[u,v])\in\frac{\chi_t}{\pi\mathsf T}\sum_{l<\mathsf T}G_l\qquad\text{for every }a\in W_i .
\end{equation}
Denote this interval by $[\underline I_t(i;j),\overline I_t(i;j)]$ when $[u,v]=W_j$.

\paragraph{The arithmetic}
All numbers are intervals whose endpoints are integer multiples of $2^{-256}$, with outward
rounding at every operation. The square root of an interval $[\ell,h]$ is enclosed by
$[\lfloor\sqrt{\max(\ell,0)\,2^{256}}\rfloor,\lceil\sqrt{\max(h,0)\,2^{256}}\rceil]/2^{256}$,
computed with exact integer square roots. The cosines at the points $l\pi/\mathsf T$ are enclosed
by the Taylor series: once the ratio of consecutive terms is at most $\frac12$, twice the first
omitted term bounds the remainder, and the summation stops when this bound is below $2^{-128}$,
the rounding widths being retained. The number $\pi$ is enclosed by
Machin's formula $\pi/4=4\arctan\frac15-\arctan\frac1{239}$ with the first omitted term of each
alternating series as remainder. No floating-point arithmetic and no elliptic integral is used.

\paragraph{The bounds}
For the upper bound, for each outer target square $(i,j)$ the program sums, over all outer source
squares $(c,d)$ (including those that cross the hypotenuse), the upper endpoints of
$8[\overline I_1(i;c)\overline I_0(j;d)+\overline I_0(i;c)\overline I_1(j;d)]$;
the maximum $U$ of these sums over $(i,j)$ is, by \eqref{eq:I-enclosure}, a bound for
$\sum_{(c,d)}P_{\mathcal R_{cd}}(x)$ for every $x$ in the target square, so
Proposition~\ref{prop:rectangles}(b) with $h\equiv1$ gives $\|B\|\le U$.
For the lower bound, use inner squares only. By \eqref{eq:I-enclosure},
$J_t(W_i,W_c)=\int_{W_i}I_t(a;W_c)\,a\dd a\ge\frac{2i+1}{2\mathsf N^2}\underline I_t(i;c)$, so
$\mathsf m_{(ij),(cd)}=8\,\frac{(2i+1)(2j+1)}{4\mathsf N^4}
[\underline I_1(i;c)\underline I_0(j;d)+\underline I_0(i;c)\underline I_1(j;d)]$ satisfies the hypothesis of
Proposition~\ref{prop:rectangles}(c), and with $\|h\|^2=1/24$ the lower endpoint of
$24\sum\mathsf m_{(ij),(cd)}$ is a lower bound $L$ for $\|B\|$. The program ends by
checking $0<L\le U$. The outputs are $L=n_L/2^{256}$ and $U=n_U/2^{256}$ with the
integers $n_L,n_U$ of Proposition~\ref{prop:enclosure}, and the decimals there are outward
roundings of $L$ and $U$ to $24$ places. This proves Proposition~\ref{prop:enclosure}; the
statement rests on Proposition~\ref{prop:rectangles}, on the correctness of the program, and on
the exactness of integer arithmetic. \qed

\bibliographystyle{elsarticle-num-names}
\bibliography{refs}
\end{document}
